\documentclass{article}%
\usepackage[T1]{fontenc}%
\usepackage[utf8]{inputenc}%
\usepackage{lmodern}%
\usepackage{textcomp}%
\usepackage{lastpage}%
\usepackage{geometry}%
\geometry{margin=2.2cm}%
%
\usepackage{amsmath}%
\usepackage{amssymb}%
\usepackage{booktabs}%
\usepackage{longtable}%
\usepackage{array}%
\usepackage{hyperref}%
\hypersetup{hidelinks}%
\title{Soluções\\[4pt]\large Comparação com o gabarito Gabarito oficial Comvest, Unicamp 2027 1ª fase, provas Q/Y\thanks{Gerado com gpt-resolve: \url{https://github.com/lgabs/gpt-resolve}}}%
\date{30/10/2026}%
%
\begin{document}%
\normalsize%
\maketitle%
\section*{Resumo}\textbf{Nota média: 10,0/10} (72 de 72 itens com a resolução correta; cada item vale o mesmo dentro da questão).\par\medskip%
\begin{longtable}{@{}c c p{0.78\textwidth}@{}}\toprule Questão & Nota & Resumo \\ \midrule
1 & 10,0 & A resolução e o gabarito oficial da Comvest marcam a alternativa B, e os dois estão certos. O trecho descreve uma relação de benefício mútuo: primeiro o apoio imperial à Igreja, depois a Igreja preservando o legado romano. Isso exclui a ruptura da (a), a autonomia sem apoio imperial da (c) e a substituição das instituições da (d). A argumentação histórica da resolução está correta. Só há ressalvas pequenas de redação, como a simplificação sobre o 'cesaropapismo' e o comentário, pertinente, sobre o anacronismo do nome 'Igreja Católica Apostólica Romana'. \\
2 & 10,0 & A resolução marca (C), a mesma resposta do gabarito oficial da Comvest, e essa resposta está correta. Pelo texto, o principal critério de classificação era selvagem/familiarizado, e não humano/não humano, e as práticas resultantes eram diferentes da pecuária e da propriedade animal europeias. A análise das alternativas erradas (a), (b) e (d) está certa. O único reparo, sem efeito no resultado, é que, na (b), a resolução acrescenta um argumento sobre o caráter 'simbólico' da incorporação que vai além do que o texto diz. \\
3 & 10,0 & A resolução marca C (voto censitário), o mesmo que o gabarito oficial da Comvest. As duas estão certas. A justificativa é sólida e historicamente correta: o art. 92 exige 100 mil réis e o art. 94 exige 200 mil réis, a eleição era indireta e só a Lei Saraiva, de 1881, excluiu os analfabetos. A resolução também descarta bem as alternativas a, b e d. Há apenas uma observação menor de redação sobre a exclusão das mulheres, que não estava expressa no texto constitucional. \\
4 & 10,0 & A resolução marca C, a mesma letra do gabarito oficial da Comvest, e as duas estão certas. A alternativa C repete o argumento central de Sollai: houve resistência regional no sul, e o Estado respondeu com estratégias de 'adaptação profunda' adaptadas ao local. A resolução descarta A, B e D com argumentos corretos, baseados no texto e no conhecimento histórico sobre a Revolução Verde do pós{-}guerra. O único reparo é de redação, numa frase da análise da alternativa (d). \\
5 & 10,0 & A resposta da resolução é (B), a mesma do gabarito oficial da Comvest, e ambos estão certos. A alternativa (b) é a única que reúne os três pontos do excerto: proteção espiritual (\textbackslash{}"fechar\textbackslash{}" o corpo), circulação ampla no Atlântico escravista e mistura de tradições africanas, europeias e americanas. A análise das alternativas erradas está correta. Há só duas imprecisões menores de redação, ambas inferências que vão além do texto e nenhuma afeta a resposta. \\
6 & 10,0 & A resolução e o gabarito oficial da Comvest coincidem: a resposta é (B), e as duas estão certas. A alternativa (b) repete a tese explícita do trecho da Comissão Nacional da Verdade, de que a sobrevivência e o território dos povos indígenas foram subordinados aos projetos desenvolvimentistas do Estado. As demais alternativas contradizem o texto, e a resolução elimina cada uma corretamente. Só há observações menores de redação e de contexto acrescentado, sem impacto na nota. \\
7 & 10,0 & A resolução e o gabarito oficial da Comvest coincidem: a resposta é (A). A análise das alternativas está certa. A (a) junta a cultura material e a memória, que vêm da imagem do museu, com a transformação da tela da TV e a ida do jogo para dentro de casa, que vêm do excerto. As outras três alternativas fazem afirmações que nenhuma das fontes sustenta. Não há nada a corrigir. \\
8 & 10,0 & A resolução e o gabarito oficial da Comvest coincidem: alternativa A (CO2, com salto de pressão seguido de queda gradual e condutividade crescente). Conferi o enunciado no PDF original e o raciocínio químico e físico está correto. A única ressalva, de baixa importância, é que a estimativa numérica de \textasciitilde{}50 $\mu$S/cm supõe p\_CO2 $\approx$ 1 atm, acima da pressão parcial real indicada pelo gráfico. Isso não afeta a resposta. \\
9 & 10,0 & A resolução está correta e coincide com o gabarito oficial (B). O critério eletroquímico está certo: só o Au$^3$? tem E$^\circ$red maior que +0,80 V, e com ele a fem é +0,13 V. Com Cu$^2$? a fem seria -0,46 V. A explicação da morfologia também está certa: 3 Ag saem para cada 1 Au depositado, o que deixa uma nanogaiola oca e porosa. As contas conferem. As únicas ressalvas são de baixa importância: o potencial dado para o Au e a simplificação do mecanismo. \\
10 & 10,0 & A resolução marca a alternativa d) e o gabarito oficial também, então as duas estão certas. Pelo gráfico, K cai com a temperatura: a reação é exotérmica e aquecer desfavorece o NH3. Como $\Delta$n\_gás = -1, aumentar a pressão favorece o NH3. O raciocínio está correto e completo, inclusive a observação de que a temperatura alta usada na indústria se explica pela cinética. Não encontrei erros. As únicas ressalvas são pequenas: o efeito da pressão não pode ser lido na figura, e a resolução supõe gás ideal sem dizer. \\
11 & 10,0 & A resolução e o gabarito oficial da Comvest concordam: a resposta é (c), A e B. A justificativa está correta. O grupo sulfônico livre torna o ligante A o mais hidrofílico, os carboxilatos de B dão afinidade intermediária e o ZIF{-}8 (C) é hidrofóbico, porque seus N estão coordenados ao Zn e o --CH3 é apolar. Só há imprecisões menores de redação, que não afetam a resposta. \\
12 & 10,0 & A resolução escolhe (C), a mesma alternativa do gabarito oficial da Comvest, e as duas estão certas. O metabolismo do metanol produz ácido fórmico, que consome o bicarbonato e baixa o pH: é uma acidose metabólica (HCO3{-} < 15 mmol/L, pH < 7,30), com indicação de tratamento quando o metanol passa de 50 mg/dL. Conferi as contas auxiliares ($\approx$ 15,6 mmol/L de metanol, pH 7,40 em condições normais, pKa $\approx$ 3,74) e estão corretas. Ficam só ressalvas menores de redação. \\
13 & 10,0 & A resolução está certa e concorda com o gabarito oficial: alternativa (D). A eletrólise da água com NaCl produz H2 no cátodo e Cl2 no ânodo. O Cl2 se desproporciona na água: Cl2 + H2O ? HClO + H+ + Cl{-}. Das quatro equações, essa é a única balanceada em átomos e em carga. A transcrição confere com o PDF, e não encontrei erro que mude a resposta. As ressalvas são só de rigor: o potencial do O2 em pH neutro e a formação de ClO{-} em meio básico. \\
14 & 10,0 & A resolução está certa e bate com o gabarito oficial: alternativa (B). No gráfico, a concentração plasmática sobe rápido depois de cada dose e se acumula, então a absorção é mais rápida que a metabolização. Os degraus da dose acumulada dão 35 mg/dia (0,7 mg/kg). O mínimo de 120 mg/kg $\times$ 50 kg = 6000 mg exige 6000/35 $\approx$ 171 dias, cerca de 170. Os 215 dias correspondem ao limite superior: 7500/35 $\approx$ 214 dias. Não há erro nas contas nem na interpretação. \\
15 & 10,0 & A resolução e o gabarito oficial concordam: alternativa (a), 95 pizzas de brigadeiro. Conferi o enunciado no PDF original e refiz as contas, inclusive a arrecadação que daria cada alternativa. Está tudo certo e não há nada a levar a um professor. \\
16 & 10,0 & A resolução está correta e coincide com o gabarito oficial: alternativa D, vermelho vence com 2 pontos. Como as distâncias são positivas, comparar os quadrados equivale a comparar as distâncias. A pedra vermelha em (1/2, 1/5) é a mais próxima do centro. A azul mais próxima tem d$^2$ = 13/4, e só duas vermelhas ficam abaixo desse valor: r$^2$ = 0,29 e r$^2$ = 2. Conferi as contas em Python com frações exatas. Não encontrei erro, ambiguidade nem hipótese escondida que precise ir a um professor. \\
17 & 10,0 & A resolução e o gabarito concordam: a resposta é a alternativa (B). O determinante é y = {-}x\^{}3+27x+10, uma cúbica com coeficiente líder negativo. Tem mínimo em ({-}3, {-}44), máximo em (3, 64), corta o eixo y em 10 e tem três raízes reais: {-}5, {-}0,37 e 5,37. Conferi numericamente em Python e comparei com o gráfico (b) da página 8 do PDF original. Não encontrei erros nem ambiguidades. \\
18 & 10,0 & A resolução e o gabarito oficial dão a mesma resposta: alternativa C (b = 10). A busca exaustiva confirma que a terna de inteiros com produto {-}8 e menor soma é ({-}1,{-}1,{-}8). A análise de casos da resolução e a verificação por expansão do polinômio estão corretas. Não há erro, só a hipótese de leitura, que a própria resolução declara. \\
19 & 10,0 & A resolução e o gabarito oficial concordam: alternativa D (07h30 e 19h30). Refiz a conta: o mínimo ocorre quando sen(($\pi$/12)(2t+3)) = -1, o que dá t = 7,5 h e t = 19,5 h no intervalo {[}0, 24{]}, com altura de 0,6 m. O enunciado transcrito confere com o original, e o raciocínio da resolução não tem erros. \\
20 & 10,0 & A resolução está correta e coincide com o gabarito oficial: alternativa C, área de B = 3/2. A área de B é a soma da série geométrica de razão 1/3, que vale 1/(1 - 1/3) = 3/2. As verificações também conferem: a área de A é 1/ln 3 $\approx$ 0,91, menor que 3/2, e as somas parciais convergem para 1,5. Não encontrei pontos duvidosos. \\
21 & 10,0 & A resolução está correta e coincide com o gabarito oficial (alternativa A, a = 1). A conta do determinante (-2r$^2$, que não depende de a), o descarte de r = 2 pela razão negativa e a verificação final estão certos. Não há divergências nem pontos duvidosos. \\
22 & 10,0 & A resolução marcou (B), igual ao gabarito oficial da Comvest, e as duas estão certas. A resolução deduz as retas dos lados, chega à dupla desigualdade |x| $\leq$ y $\leq$ 2 - |x| e mostra que daí já sai |x| $\leq$ 1, sem precisar impor isso à parte. Também confere o resultado pela forma |x| + |y - 1| $\leq$ 1 e descarta (a), (c) e (d) com contraexemplos válidos. Não encontrei erro nem ambiguidade relevante. \\
23 & 10,0 & A resolução e o gabarito dão a mesma resposta: alternativa C, {[}1,2{]}. O argumento da resolução está correto e completo. A parábola tem vértice em x = 1/2. Os intervalos {[}0,2{]}, {[}-1,2{]} e {[}-2,1{]} contêm pares simétricos em torno de 1/2, como f(0) = f(1) e f(-1) = f(2), então f não é injetora neles. Em {[}1,2{]}, f é estritamente crescente, e a inversa está bem calculada. \\
24 & 10,0 & A resolução e o gabarito oficial concordam: alternativa (D), 149/198. O argumento está correto: o mdc só é par se os dois números forem pares, o que dá P = 1 - C(50,2)/C(100,2). Uma enumeração exaustiva em Python confirma o valor (3725/4950 = 149/198). Não há pontos a levar ao professor. \\
25 & 10,0 & A resolução está correta e bate com o gabarito oficial (D). A área é A = (1/4)(1/2 - f(1/2)). Como 1/8 < f(1/2) < 1/4, segue que 1/16 < A < 3/32. A transcrição confere com o original, e as contas foram checadas com frações exatas. Não há ponto duvidoso. \\
26 & 10,0 & A resolução está correta e coincide com o gabarito oficial (alternativa A). Os triângulos OCB e OAC são isósceles, o que dá $\angle$AOP = 45$^\circ$ pelo ângulo externo. Pela lei dos cossenos, AP = $\surd$(2 - $\surd$2) $\approx$ 0,765. A geometria foi confirmada numericamente. O único deslize são dois valores aproximados das alternativas descartadas (c e d), sem efeito na resposta. \\
27 & 10,0 & A resolução e o gabarito oficial da Comvest marcam a mesma alternativa, D, e ambos estão certos. \textbackslash{}"Fantasia\textbackslash{}" aparece na acepção 4 (traje), em sentido figurado: o folião \textbackslash{}"veste\textbackslash{}" o imperativo moral \textbackslash{}"Não mate sua mina\textbackslash{}". A transcrição bate com o original. A justificativa da resolução está correta e descarta cada alternativa errada associando{-}a às acepções 1, 2 e 3. Nada a levar ao professor. \\
28 & 10,0 & A resposta da resolução (B) é a mesma do gabarito oficial da Comvest. A interpretação está bem apoiada no texto: o desengano vem do fim da festa (\textbackslash{}"Quarta{-}feira sempre desce o pano\textbackslash{}"), quando o eu lírico volta à realidade. As alternativas erradas são descartadas com trechos da própria canção. Não encontrei nenhum ponto que mereça revisão. \\
29 & 10,0 & A resolução e o gabarito oficial da Comvest coincidem: a resposta é a alternativa D (presente do subjuntivo: viva, possa e saiba). A justificativa está correta. Ela relaciona a expectativa agostiniana ao subjuntivo optativo da estrofe introduzida por \textbackslash{}"esperança\textbackslash{}" e descarta bem os outros tempos verbais, ligados à memória ou à atenção. Não encontrei erro, e restam só observações menores sobre a leitura dos distratores. \\
30 & 10,0 & A resolução marca A, a mesma letra do gabarito oficial da Comvest, e as duas estão certas. A transcrição bate com o original. A resolução justifica a resposta com trechos do próprio texto (\textbackslash{}"rebaixamento\textbackslash{}", \textbackslash{}"etiqueta de gênero\textbackslash{}", \textbackslash{}"preconceito com açúcar\textbackslash{}") e descarta bem as outras alternativas. A B é a ofensa explícita (\textbackslash{}"veado\textbackslash{}"). A C fala de empatia, que o texto não menciona. A D descreve a reação de quem recebe o elogio, e o texto liga a resposta defensiva ao primeiro insulto. Não encontrei nada que precise ir a um professor. \\
31 & 10,0 & A resolução marca a alternativa C, a mesma do gabarito oficial da Comvest, e as duas estão certas. Os dois textos chegam a conclusões opostas, mas ambos usam dados estatísticos para sustentar sua tese. As alternativas a, b e d foram descartadas pelos motivos certos. A única falha é um erro de grafia no nome do autor do Texto 1, que não muda a resposta. \\
32 & 10,0 & A resolução e o gabarito oficial concordam: a resposta é a alternativa B. Pelo texto 2, com controles institucionais, os mesmos dados mostram que o Brasil trabalha mais do que o esperado, e a hipótese cultural é rejeitada. Isso reverte a conclusão do texto 1. A justificativa da resolução é coerente com o enunciado original e descarta corretamente as alternativas A, C e D. Não encontrei pontos que precisem ir a um professor. \\
33 & 10,0 & A resolução e o gabarito oficial da Comvest coincidem: a resposta da questão 33 é a C. A interpretação está correta e bem fundamentada no texto, com três marcas da relação senhor--escravizado: animalização, violência e silenciamento. Ela se apoia também no capítulo \textbackslash{}"O vergalho\textbackslash{}", citado com precisão, e descarta as outras alternativas com bons argumentos. Os únicos problemas são de forma: um erro de LaTeX (\textbackslash{}"\textbackslash{}\textbackslash{}emph\{reproduz*\textbackslash{}") e um trecho em que narrador e autor poderiam estar mais bem separados. \\
34 & 10,0 & A resolução e o gabarito oficial da Comvest marcam a mesma alternativa, D, e ela está correta. A advertência faz parte da ficção da obra: Lima Barreto finge ser apenas o editor e atribui a autoria ao narrador{-}personagem Augusto Machado. A justificativa da resolução está certa e não traz imprecisões que precisem ir a um professor. \\
35 & 10,0 & A resolução está certa e coincide com o gabarito oficial (C). A bengala, lida como o cajado do pastor, faz do eu{-}lírico o condutor dos próprios passos, o que expressa o protagonismo de quem conduz. As alternativas erradas foram descartadas com argumentos tirados do texto. Só há duas imprecisões pequenas no contexto (o título do livro e a caracterização da ode), que não afetam a resposta. \\
36 & 10,0 & A resolução e o gabarito oficial da Comvest marcam a mesma alternativa, (D), e ela está correta. Na fala da autora, a dor faz parte do processo de criação, por isso (c) está errada. Já (a) e (b) atribuem à escrita um efeito de \textbackslash{}"diminuir\textbackslash{}" o abandono ou o silêncio, o que não acontece. A transcrição confere com o original, e a argumentação está bem apoiada nos excertos. Só há ressalvas pequenas de redação, que não mudam a resposta. \\
37 & 10,0 & A resolução e o gabarito oficial da Comvest coincidem: a resposta é (B). A interpretação está correta, porque a troca de Grace por Afamefuna, o nome igbo dado pela avó, representa a volta às raízes anteriores à colonização. As alternativas (a), (c) e (d) foram descartadas com justificativas adequadas. Não há erro a apontar, só uma observação de redação sem efeito na nota. \\
38 & 10,0 & A resolução e o gabarito oficial da Comvest marcam a alternativa C, e as duas estão certas. Pelo contexto do conto, a esposa de Montiel interpreta a exploração como caridade, e o narrador expõe essa ignorância ao leitor. Isso é ironia dramática construída pelo narrador. A justificativa da resolução está correta e descarta bem as outras alternativas: (a) é antítese, não paradoxo; em (b) o raciocínio é coerente; em (d) a esposa não sabe das intenções do marido. Há só uma imprecisão de nomenclatura sem importância. \\
39 & 10,0 & A resolução está certa: marca (B), igual ao gabarito oficial da Comvest. A justificativa segue o Texto 2, que descarta explicações demográficas, censitárias e de cálculo por benefício e atribui a mudança às ideologias raciais que convivem entre si. As leituras do gráfico também estão corretas. Há apenas observações menores de interpretação, que não afetam a nota. \\
40 & 10,0 & A resolução e o gabarito oficial da Comvest marcam a mesma alternativa, D, e os dois estão corretos. É uma questão de interpretação de texto. A resolução identifica bem a tese do autor: em contextos de polarização, a democracia deve ser defendida como método de administrar os conflitos em paz e sem repressão. Ela também descarta corretamente A, B e C, que falam em crise, em 'maioria' ou na impossibilidade de administrar os conflitos sem repressão, ideias ausentes do texto ou contrárias a ele. \\
41 & 10,0 & A resolução e o gabarito oficial da Comvest marcam a mesma alternativa, (C), e ela está correta. A transcrição do enunciado confere com a prova original. A alternativa (c) reproduz o mecanismo descrito no Texto 1: os negacionistas se apropriam de retóricas e credenciais científicas, fabricam falsas controvérsias e com isso criam dissensos e incertezas. Também corresponde à interseção sombreada do Texto 2. Os descartes de (a), (b) e (d) estão bem justificados com trechos do texto. Não há erro. \\
42 & 10,0 & A resolução está certa e coincide com o gabarito oficial: alternativa (A). A tirinha mostra a transcrição. O DNA serve de molde, e o RNAm é complementar e tem U no lugar de T (CACA CTCT ACT ACT $\to$ GUGU GAGA UGA UGA, pareamento conferido). O processo ocorre no núcleo, e o RNAm leva a informação ao citoplasma. As alternativas (b), (c) e (d) foram descartadas com justificativas corretas. Não há divergência. \\
43 & 10,0 & A resolução marca (C), a mesma resposta do gabarito oficial da Comvest, e as duas estão certas. A (c) descreve a seleção positiva citada no texto: a variante duplicada é vantajosa e fica mais comum ao longo das gerações. As outras alternativas confundem fluxo gênico, deriva e mutação ou tratam a especiação como súbita. Há só uma imprecisão de redação na análise da (b), que não afeta a resposta. \\
44 & 10,0 & A resolução está certa e coincide com o gabarito oficial: alternativa B. As sequências são cnidários $\to$ platelmintos, membrana plasmática e entrada de Na+ $\to$ saída de K+. A justificativa está correta em todas as lacunas e descarta bem as outras alternativas. Só há ressalvas pequenas de redação, que não mudam a resposta. \\
45 & 10,0 & A resolução e o gabarito oficial da Comvest marcam a mesma alternativa, a (C), e os dois estão certos. Pela figura, a seta 2 está no lado sombreado do caule e a seta 3 no lado iluminado. A auxina se concentra no lado sombreado, que se alonga mais, e isso produz o fototropismo positivo. As alternativas (b) e (d) erram ao atribuir o gravitropismo à citocinina, e a (a) põe o aumento de auxina no lado iluminado. Só há ressalvas de redação, sem efeito na nota. \\
46 & 10,0 & A resolução e o gabarito oficial dão a mesma resposta: alternativa B. O raciocínio está correto. Pirâmide de energia não se inverte, o que elimina a e c, então são pirâmides de biomassa. Em t2 produtores e consumidores primários caíram e os níveis superiores ainda não, o que mostra uma teia em transição e não um equilíbrio restabelecido, eliminando a d. A leitura da figura confere com o original. Não há erro relevante, só uma imprecisão de redação: a resolução fala em 'conservação da energia' onde caberia falar em perda de energia como calor. \\
47 & 10,0 & A resolução e o gabarito oficial da Comvest coincidem na alternativa D, e a resposta está correta. A AIDS é o estágio mais grave da infecção pelo HIV, e a terapia antirretroviral, ao deixar a carga viral indetectável, impede a progressão para a doença. Os erros apontados pela resolução nas outras alternativas também estão certos: em A, a ideia de eliminação natural do vírus; em B, a negação de que o tratamento previne a transmissão; em C, a relação causal invertida entre AIDS e doenças oportunistas. A transcrição bate com o original e não há pontos a levar a um professor. \\
48 & 10,0 & A resolução e o gabarito oficial da Comvest dão a mesma resposta, (C), e ela está correta. Na medula espinhal, os axônios se concentram na substância branca. Ela fica na periferia e é formada pelos tratos ascendentes e descendentes, enquanto a substância cinzenta central reúne os corpos celulares. A justificativa da resolução e a análise das alternativas estão corretas e completas. Não achei erro nem no gabarito nem na resolução. \\
49 & 10,0 & A resolução e o gabarito oficial da Comvest dão a mesma resposta, a alternativa D, e ela está correta. Pelo texto, as três dimensões de Milton Santos (fábula, perversidade e outra globalização como tendência) formam uma unidade. O mundo aparece como uma totalidade aberta e em movimento, com transformações já em curso nas grandes cidades. Os argumentos usados para descartar a, b e c estão de acordo com o texto e com a obra do autor. Não encontrei nenhum ponto que precise ser levado a um professor. \\
50 & 10,0 & A resolução e o gabarito oficial da Comvest coincidem: alternativa B. Conferi o enunciado no PDF original e a transcrição está fiel. A argumentação da resolução está correta. Ela descarta a unipolaridade (a), a idealização ética da China (c) e a bipolaridade com economias 'independentes' e 'sem critérios éticos' (d). A única ressalva é leve: a expressão 'hegemonia num mundo multipolar' da própria alternativa B tem uma pequena tensão conceitual, mas isso não muda a resposta. \\
51 & 10,0 & A resolução marca (A), igual ao gabarito oficial da Comvest, e está correta. O texto descreve as usinas flex, que processam cana e milho. As outras alternativas foram descartadas pelos motivos certos: a cana é perecível e o milho pode ser estocado (b); a cana é semiperene e o milho vem da safrinha depois da soja (c); a colheita de milho é mecanizada (d). Só há imprecisões menores de redação, que não afetam a nota. \\
52 & 10,0 & A resolução e o gabarito oficial concordam: a resposta é a alternativa (A). Os números da tabela estão certos, e os municípios foram associados corretamente aos vetores de crescimento: as fronteiras agrícola, mineral, do petróleo e do turismo, de um lado, e a dispersão metropolitana com expansão imobiliária, de outro. As alternativas erradas foram descartadas pelo motivo certo. Não encontrei nenhum erro. \\
53 & 10,0 & A resolução e o gabarito oficial concordam: a resposta é (a). Conferi pelo mapa original. O deserto tem um único rótulo A, que cobre a maior parte do bioma, e rótulos C só nas bordas. A savana aparece nas três zonas, A, B e C. A (b) inverte os biomas, a (c) contraria o mapa e a (d) também, porque há C sobre a savana. A argumentação da resolução está correta. Algumas leituras finas de rótulos com linha de chamada podem ser discutíveis, mas não mudam a alternativa. \\
54 & 10,0 & A resolução e o gabarito oficial concordam: alternativa (C). Conferi os mapas no PDF original ampliado, e a leitura de classes da resolução está certa. Na RM de São Paulo, P fica em 1600--1800 e a ETP abaixo de 1600, então P > ETP sem ambiguidade. A justificativa também está certa: sazonalidade, armazenamento e demanda exigem monitoramento e alertas. As eliminações de (a), (b) e (d) estão bem fundamentadas. Só a transcrição diverge do original na precipitação de Vitória, e isso não altera a resposta. \\
55 & 10,0 & A resolução está certa e bate com o gabarito oficial Comvest: a resposta é a alternativa A. Ela mostra que o logotipo é um símbolo institucional inspirado no traçado radial e concêntrico do campus. Como não tem os elementos de um mapa (escala, orientação, legenda e correspondência de posições), não é uma representação cartográfica, o que elimina as alternativas b, c e d. Não há divergência a resolver. \\
56 & 10,0 & A resolução e o gabarito oficial da Comvest dão a mesma resposta, a alternativa B, e as duas estão certas. A leitura da figura confere com o PDF original. O ponto central da resolução também está correto: as barras inferiores valem só para o subgrupo de 51\% que apoia 'some' (algumas deportações). Por isso, para chegar ao total de entrevistados, é preciso multiplicar por 0,51, e é isso que derruba (c) ($\approx$26,5\%) e (d) ($\approx$81,5\%). Há só uma imprecisão de redação, o 'no máximo' em (d), que não afeta a resposta. \\
57 & 10,0 & A resolução e o gabarito oficial da Comvest marcam a alternativa A, e os dois estão certos. Moura chama a arte de 'revolutionary' e Morrison, de 'dangerous': os dois veem a arte como transgressora. Só Morrison diz quem ataca os artistas (ditadores, ocupantes de cargos e quem quer controlar e enganar). As outras alternativas foram descartadas pelos motivos certos: a B inverte o que Moura diz, a C usa o falso cognato deceive/decepção e a D afirma algo que nenhum dos textos diz. Não há pontos de dúvida. \\
58 & 10,0 & A resolução marca (C), igual ao gabarito oficial da Comvest. Pela ilustração do PDF original, (C) está correta: são as cápsulas vermelhas, em forma de gota e cobertas de cerdas, do urucum, que é o fruto descrito no texto. A análise das alternativas e a identificação pela etimologia tupi estão corretas, e não encontrei pontos duvidosos. \\
59 & 10,0 & A resolução e o gabarito oficial da Comvest escolhem a mesma alternativa, a (B), e ela está certa. O texto apresenta as gírias como marca de pertencimento ao grupo (in{-}group status) e como instrumento para construir comunidade e identidade compartilhada. As alternativas a, c e d são descartadas com trechos do próprio texto. Há só um asterisco sobrando na resolução, sem efeito na resposta. \\
60 & 10,0 & A resolução e o gabarito oficial da Comvest marcam a mesma alternativa, (C), e os dois estão certos. A alternativa (c) traduz quase literalmente a frase \textbackslash{}"can commodify intimacy through emotionally manipulative design\textbackslash{}". A resolução explica bem por que (a), (b) e (d) contradizem o texto ou trocam o papel dos conceitos. Não há erro a apontar. \\
61 & 10,0 & A resolução e o gabarito oficial da Comvest marcam a mesma alternativa, (D), e os dois estão certos. O poema usa passado e particípio para os traumas coloniais (stripped, erased, buried, burned, bled) e presente simples para a resistência (breathes, births, dances, wear, rise). A justificativa da resolução está correta, inclusive ao descartar (a): quem simboliza o sofrimento é a cruz, e a coroa simboliza dignidade. Os pontos levantados são de baixa importância: o termo \textbackslash{}"infringido\textbackslash{}" no lugar de \textbackslash{}"infligido\textbackslash{}" na prova e alguns particípios com sentido positivo (crowned). \\
62 & 10,0 & A resolução e o gabarito oficial da Comvest coincidem: a resposta é a alternativa D, e ela está correta. Pelo texto, todos os seres sencientes, humanos ou não, têm interesse em não sofrer e em não morrer, e o uso de animais como propriedade fere esses dois interesses. A resolução refuta bem a (a), que contradiz a afirmação de que os animais explorados são sencientes; a (b), já que o texto não fala em criminalizar e diz que a exploração é injustificável mesmo quando \textbackslash{}"humane\textbackslash{}"; e a (c), já que julgar pessoas \textbackslash{}"is never a good idea\textbackslash{}". Há apenas uma pequena imprecisão de redação ao citar a alternativa (a), que não afeta o resultado. \\
63 & 10,0 & A resolução e o gabarito oficial concordam na alternativa (a). Refiz a conta: d = c$\cdot$$\Delta$t = 3,0$\times$10\^{}8 $\times$ (1,3$\times$10\^{}9 $\times$ 3,2$\times$10\^{}7) $\approx$ 1,25$\times$10\^{}25 m $\approx$ 1,2$\times$10\^{}25 m. O enunciado transcrito confere com o PDF. A análise dos distratores (esquecer a conversão de anos para segundos ou tomar bilhão como 10\^{}6) está correta. Não há erro a apontar. \\
64 & 10,0 & A resolução e o gabarito oficial da Comvest marcam a mesma alternativa, B, e ambos estão certos. Como g = GM/R$^2$, a aceleração da gravidade não depende da massa m do corpo, é diretamente proporcional à massa da Terra e inversamente proporcional ao quadrado do raio. A dedução da resolução está correta, e as verificações de unidade e de ordem de grandeza ($\approx$ 9,8 m/s$^2$) também. Não há pontos para levar a um professor. \\
65 & 10,0 & A resolução e o gabarito oficial dão a mesma resposta, a alternativa (d), 240 N. A física está correta: a carga pende de um dos ramos do fio que passa pela polia móvel, e o fio da manopla puxa o eixo dessa polia. Por isso F = 2mg = 240 N; a montagem dobra a força em vez de reduzi{-}la à metade. Ampliei a figura original a 600 dpi e ela confirma essa leitura. O único reparo é de redação: a resolução fala em uma normal do piso, mas na figura o bloco já está suspenso. Isso não altera o resultado. \\
66 & 10,0 & A resolução e o gabarito concordam: alternativa C, 4,0 N. Durante o deslizamento o atrito é cinético, e f\_at = $\mu$\_C$\cdot$m$\cdot$g = 0,02$\cdot$20$\cdot$10 = 4,0 N. A física, a conta e a explicação dos distratores estão certas, e a transcrição confere com o original. \\
67 & 10,0 & A resolução e o gabarito oficial concordam: a resposta é (B). Pela calorimetria, $\Delta$T = Q/(m$\cdot$c) = 80/(10$\cdot$0,4) = 20 $^\circ$C, então T3 = 50 $^\circ$C. O gráfico (b) é o único com aquecimento do sólido, patamar de fusão a 30 $^\circ$C e aquecimento do líquido até 50 $^\circ$C em $\approx$ 230 s. A resolução elimina (a), (c) e (d) corretamente e não tem erro de conta nem de física. \\
68 & 10,0 & A resolução e o gabarito oficial da Comvest marcam a mesma alternativa, B (0,25 A), e ambos estão corretos: i = P/U = 3,0 W / 12 V = 0,25 A. Conferi o enunciado no PDF original e ele coincide com a transcrição. Não encontrei erro nem ambiguidade a apontar. \\
69 & 10,0 & A resolução está correta e bate com o gabarito oficial (D). Com a chave fechada, os quatro resistores de 1,2 k$\Omega$ ficam em paralelo, dando R\_eq = 300 $\Omega$ e i = 9,0 V/300 $\Omega$ = 30 mA. A justificativa de cada distrator também está certa. Não há divergência. \\
70 & 10,0 & A resolução e o gabarito oficial da Comvest coincidem: a resposta é a alternativa D. A transcrição confere com o original. A leitura do trecho de Habermas e a eliminação das alternativas a, b e c estão corretas e bem fundamentadas no texto. Não encontrei nenhum ponto que mereça ser levado a um professor. \\
71 & 10,0 & Questão 71: a resolução marca C e o gabarito oficial da Comvest também dá C, então as duas estão certas. Para Weil, a força reifica quem a sofre e também quem a usa. As alternativas A, B e D foram descartadas com argumentos corretos tirados dos textos. O único ponto anotado é de redação e não muda a resposta. \\
72 & 10,0 & A resolução marca (A), e o gabarito oficial também. Conferi o enunciado no PDF e as duas respostas estão certas: o texto diz que o ego moderno se constitui também como conquistador e colonizador do Outro, o que vai além do cogito cartesiano. A resolução descarta corretamente (b), (c) e (d) com base no próprio texto. Não há erro a apontar, só uma observação de pouca importância sobre o uso de um conhecimento que não está no texto. \\
\bottomrule\end{longtable}%
\section*{Questão 1}%
\begin{longtable}{@{}p{0.06\textwidth} p{0.33\textwidth} p{0.33\textwidth} p{0.1\textwidth} p{0.08\textwidth}@{}}\toprule Item & Resolução & Gabarito & Veredito & Certo \\ \midrule
1 & (B): relação por vezes simbiótica entre Igreja e Estado romano, que fortaleceu as duas instituições e favoreceu determinados grupos sociais da Antiguidade Tardia. & B & igual & ambos \\
\bottomrule\end{longtable}%
\begin{itemize}%
\item%
\textbf{1)} A transcrição confere com o original (p. 3 do PDF). O trecho descreve duas fases. Primeiro, a Igreja se beneficiou do apoio imperial. Depois, ajudou a preservar o legado de Roma no século V. Isso é reforço mútuo, que corresponde à alternativa (b). A (a) fala em ruptura ou separação, o que contradiz o texto e é anacrônico. A (c) põe a expansão depois do colapso e sem apoio imperial, o que também contradiz o trecho. A (d) troca 'preservar' por 'substituir', distorcendo o que o texto diz. Os dados históricos usados na resolução estão corretos: Galério (311), Milão (313), Niceia (325), Tessalônica (380), Rômulo Augústulo (476), Calcedônia convocada por Marciano (451), exílio de Atanásio sob Constâncio II, Juliano (361--363), penitência de Teodósio imposta por Ambrósio (390), Sidônio Apolinário, episcopalis audientia e Cisma de 1054.%
\end{itemize}%
\paragraph{Pontos de dúvida}%
\begin{itemize}%
\item%
\textit{(statement)} A alternativa (b) usa 'Igreja Católica Apostólica Romana' para a Antiguidade Tardia, um termo anacrônico. A resolução aponta isso e trata como imprecisão de nomenclatura, que não invalida a alternativa. A leitura é razoável e é a mesma do gabarito, mas um aluno mais rigoroso poderia hesitar por causa disso.%
\item%
\textit{(solution)} A resolução diz que Calcedônia (451) é 'base do chamado cesaropapismo bizantino'. É uma simplificação: o conceito é historiograficamente debatido e costuma ser associado sobretudo a Justiniano. Não afeta a resposta.%
\end{itemize}%
\section*{Questão 2}%
\begin{longtable}{@{}p{0.06\textwidth} p{0.33\textwidth} p{0.33\textwidth} p{0.1\textwidth} p{0.08\textwidth}@{}}\toprule Item & Resolução & Gabarito & Veredito & Certo \\ \midrule
Q2 & (C) & C & igual & ambos \\
\bottomrule\end{longtable}%
\begin{itemize}%
\item%
\textbf{Q2)} Conferi o enunciado na página 3 do PDF e a transcrição está fiel. O texto diz que muitos povos ameríndios distinguiam 'sobretudo' seres selvagens e seres familiarizados, 'em vez de separar rigidamente humanos e não humanos'. Diz também que incorporavam os animais europeus segundo referências 'diferentes das concepções europeias de pecuária e propriedade animal'. A alternativa (c) resume as duas ideias. As outras três têm um erro cada: a (a) chama de familiarizados os seres que 'não faziam parte do convívio cotidiano', quando o texto diz que eram justamente os integrados a ele; a (b) põe a distinção humano/não humano como base da relação, o oposto do texto; a (d) aproxima essas práticas das concepções europeias, o que o texto nega. A resolução descarta cada alternativa pelo motivo certo.%
\end{itemize}%
\paragraph{Pontos de dúvida}%
\begin{itemize}%
\item%
\textit{(solution)} Na alternativa (b), a resolução também diz que a incorporação dos animais era 'concreta' e não apenas 'simbólica'. O texto não trata disso de forma explícita, então esse é um argumento a mais e discutível. Isso não muda nada, porque o erro decisivo da (b) é pôr a distinção humano/não humano como base da relação, e a resolução aponta esse erro corretamente.%
\end{itemize}%
\section*{Questão 3}%
\begin{longtable}{@{}p{0.06\textwidth} p{0.33\textwidth} p{0.33\textwidth} p{0.1\textwidth} p{0.08\textwidth}@{}}\toprule Item & Resolução & Gabarito & Veredito & Certo \\ \midrule
Q3 & C (voto censitário: renda líquida anual mínima de 200 mil réis, além da exclusão dos libertos e dos criminosos pronunciados) & C & igual & ambos \\
\bottomrule\end{longtable}%
\begin{itemize}%
\item%
\textbf{Q3)} A transcrição confere com o original (p. 3 do PDF). O art. 94, inciso I, condiciona o direito de ser eleitor a uma renda líquida anual de 200 mil réis 'por bens de raiz, indústria, comércio, ou emprego'. Isso é o voto censitário, e os incisos II e III são as 'outras restrições legais' citadas na alternativa C. As demais estão erradas. A) não pode ser, porque o próprio trecho exclui grupos. B) não pode ser, porque a Constituição de 1824 não exigia instrução, e os analfabetos só foram excluídos pela Lei Saraiva (1881). D) também não, porque o texto não usa cor ou origem como critério. A exclusão dos libertos se baseia na condição jurídica de ex{-}escravizado. Os dados históricos da resolução estão certos: o art. 92 exige 100 mil réis para votar nas assembleias paroquiais, o art. 94 exige 200 mil réis para ser eleitor, e a eleição era indireta, em dois graus.%
\end{itemize}%
\paragraph{Pontos de dúvida}%
\begin{itemize}%
\item%
\textit{(solution)} Na alternativa a), a resolução diz que mulheres e escravizados não tinham direitos políticos. Isso é historicamente correto, mas a Constituição de 1824 não excluía as mulheres de forma expressa. A exclusão vinha do costume e da interpretação da época. Também não aparece no trecho do enunciado. É uma imprecisão de redação e não muda a resposta.%
\item%
\textit{(statement)} A alternativa d) pode atrair candidatos, porque a exclusão dos libertos atingia, na prática, a população negra. Mesmo assim, o critério legal explícito era a condição jurídica, e não a cor, e a alternativa C é a que melhor descreve o critério central do texto, que é a renda. Não há ambiguidade real quanto à resposta.%
\end{itemize}%
\section*{Questão 4}%
\begin{longtable}{@{}p{0.06\textwidth} p{0.33\textwidth} p{0.33\textwidth} p{0.1\textwidth} p{0.08\textwidth}@{}}\toprule Item & Resolução & Gabarito & Veredito & Certo \\ \midrule
Questão 4 (única) & C & C & igual & ambos \\
\bottomrule\end{longtable}%
\begin{itemize}%
\item%
\textbf{Questão 4 (única))} A transcrição confere com o original (págs. 3--4 do PDF). Segundo o texto, a Batalha do Trigo funcionou no norte e fracassou no sul semiárido, onde os agricultores resistiram. Essa resistência levou o Estado a financiar estratégias de 'adaptação profunda', voltadas à agrobiodiversidade local, e é exatamente isso que a alternativa C diz. A alternativa A contradiz o texto, que afirma que a versão fascista 'não foi um bloco homogêneo'. A B está errada porque a Revolução Verde do pós{-}guerra se apoiou justamente em pesquisa científica e melhoramento genético. A D inverte a causa: a resistência veio da tentativa de impor as novas variedades, não de um abandono do sul pelo Estado. A resolução analisa cada alternativa corretamente e chega à mesma letra do gabarito.%
\end{itemize}%
\paragraph{Pontos de dúvida}%
\begin{itemize}%
\item%
\textit{(solution)} Na análise da alternativa (d), a frase 'O Estado não abandonou o sul, e foi isso que provocou a resistência' é ambígua. Pode ser lida como se a resistência tivesse vindo da 'presença' do Estado, quando o que a provocou foi a imposição das novas variedades. A explicação logo em seguida esclarece o sentido, então é só uma questão de redação, sem efeito na resposta.%
\end{itemize}%
\section*{Questão 5}%
\begin{longtable}{@{}p{0.06\textwidth} p{0.33\textwidth} p{0.33\textwidth} p{0.1\textwidth} p{0.08\textwidth}@{}}\toprule Item & Resolução & Gabarito & Veredito & Certo \\ \midrule
Q5 & (B) As bolsas de mandinga eram recursos espirituais de proteção usados em todo o Atlântico escravista, combinando elementos africanos, europeus (católicos) e americanos em novas práticas religiosas. & B & igual & ambos \\
\bottomrule\end{longtable}%
\begin{itemize}%
\item%
\textbf{Q5)} Conferi a transcrição com o original (p. 4 do PDF, p. 2 da prova) e ela está fiel, assim como as alternativas. A alternativa (b) segue diretamente o excerto: as bolsas serviam para "fechar" o corpo, ou seja, eram proteção espiritual. Circulavam por redes entre Minas, Recife, Rio, Lisboa, Costa da Mina e Angola, o que mostra o uso amplo no Atlântico escravista. E reuniam elementos "fruto do contato entre diferentes tradições culturais", o que é a articulação de tradições. A (a) contradiz o texto, que fala em hóstias, pedra d'ara e flora e fauna americanas. A (c) fala em "substituição", mas o texto só diz que os elementos foram incorporados. A (d) introduz uma matriz islâmica e a leitura do Alcorão, que não aparecem no excerto, e reduz as bolsas a uma única tradição. A resolução descarta cada alternativa errada com o motivo certo.%
\end{itemize}%
\paragraph{Pontos de dúvida}%
\begin{itemize}%
\item%
\textit{(solution)} Na alternativa (d), a resolução diz que "os escritos citados são orações e símbolos cristãos". O excerto fala apenas em "orações, cruzes e outros símbolos gráficos" e não diz que todos são cristãos. Além disso, a historiografia reconhece uma raiz islâmica (mandinga/malê) nesses amuletos. Isso não muda a resposta: (d) continua errada por atribuir às bolsas uma matriz exclusivamente islâmica e uma função de leitura do Alcorão que o texto não sustenta. Mesmo assim, a frase da resolução vai além do que o excerto diz.%
\item%
\textit{(solution)} No Passo 1, a resolução diz que as partes do corpo humano estavam "ligadas a práticas mágicas de matriz africana". O excerto não faz essa atribuição, então é uma inferência da resolução. Não afeta a conclusão.%
\end{itemize}%
\section*{Questão 6}%
\begin{longtable}{@{}p{0.06\textwidth} p{0.33\textwidth} p{0.33\textwidth} p{0.1\textwidth} p{0.08\textwidth}@{}}\toprule Item & Resolução & Gabarito & Veredito & Certo \\ \midrule
Questão 6 (alternativa única) & (B) & B & igual & ambos \\
\bottomrule\end{longtable}%
\begin{itemize}%
\item%
\textbf{Questão 6 (alternativa única))} Conferi a transcrição com a página 4 do PDF e ela está fiel. A alternativa (b) parafraseia diretamente a frase central do texto: "a sobrevivência dos povos indígenas e a existência de suas culturas foram subordinadas aos grandes projetos de colonização do país". O exemplo da BR{-}80 e a fala do presidente da Funai reforçam essa leitura. (a) contradiz a ideia de subordinação. (c) contradiz a falta de controle social e o caráter autoritário do regime, e é o próprio órgão indigenista que justifica a obra. (d) contradiz o caso da BR{-}80, que cortou o Parque do Xingu, dividiu os Tukarramãe e levou doenças a eles. A análise da resolução está correta em todas as alternativas.%
\end{itemize}%
\paragraph{Pontos de dúvida}%
\begin{itemize}%
\item%
\textit{(solution)} A resolução acrescenta contexto que não está no trecho: o Plano de Integração Nacional (1970), a Transamazônica e o lema "integrar para não entregar". O contexto é historicamente correto e não afeta a resposta, mas não vem do texto da prova.%
\item%
\textit{(solution)} Na análise de (d), a resolução diz que as rodovias "agravaram os conflitos com colonos e fazendeiros". O texto cita as arbitrariedades de colonos e fazendeiros, mas não diz que as rodovias agravaram esses conflitos. É uma inferência razoável, não uma afirmação literal do trecho, e não muda a eliminação da alternativa.%
\end{itemize}%
\section*{Questão 7}%
\begin{longtable}{@{}p{0.06\textwidth} p{0.33\textwidth} p{0.33\textwidth} p{0.1\textwidth} p{0.08\textwidth}@{}}\toprule Item & Resolução & Gabarito & Veredito & Certo \\ \midrule
Q7 & (A) & A & igual & ambos \\
\bottomrule\end{longtable}%
\begin{itemize}%
\item%
\textbf{Q7)} Conferi a transcrição com as páginas originais do PDF (figura da exposição Game Masters, excerto de Carreiro e alternativas a--d), e está fiel. A alternativa (a) é a única que se apoia nas duas fontes. A imagem mostra uma cultura material musealizada, ligada à memória. O excerto mostra que a tela da TV deixou de ser 'expositiva' e passou a ser 'interativa', e que o jogo saiu dos espaços públicos para dentro de casa, o que muda a sociabilidade e o consumo. A (b) fala em mudança no modo de filmar e exibir do cinema e da TV, a (c) em adaptação às técnicas de Hollywood e aprofundamento das diferenças entre linguagens, e a (d) em recusa da TV e perda de espaço; nenhuma tem apoio nas fontes. Na (d), o console até reforça a centralidade do televisor. A justificativa da resolução está correta e segue a lógica do enunciado.%
\end{itemize}%
\section*{Questão 8}%
\begin{longtable}{@{}p{0.06\textwidth} p{0.33\textwidth} p{0.33\textwidth} p{0.1\textwidth} p{0.08\textwidth}@{}}\toprule Item & Resolução & Gabarito & Veredito & Certo \\ \midrule
Q8 & (A) CO2; a pressão dá um salto logo após a liberação e depois cai aos poucos, e a condutividade cresce de forma contínua até \textasciitilde{}60 $\mu$S/cm & A & igual & ambos \\
\bottomrule\end{longtable}%
\begin{itemize}%
\item%
\textbf{Q8)} A transcrição confere com o original (p. 5 do PDF). Os gráficos de a) e c) são iguais, e os de b) e d) também. O raciocínio da resolução está correto em três pontos. (1) Ao entrar gás no frasco fechado, a pressão sobe obrigatoriamente, para qualquer gás. Por isso os gráficos com pressão constante (b, d) ficam descartados. (2) O He é inerte e quase insolúvel, então não geraria íons nem faria a pressão cair. Isso descarta c (e d). (3) O CO2 se dissolve e forma H+ e HCO3{-}: a pressão cai aos poucos e a condutividade cresce de forma monotônica, sem salto seguido de queda. Refiz no Python a estimativa de ordem de grandeza: {[}H+{]} = $\surd$(4,5e{-}7$\cdot$3,4e{-}2$\cdot$1) $\approx$ 1,24e{-}4 mol/L (pH $\approx$ 3,9) e $\kappa$ $\approx$ 395$\cdot$1,24e{-}7 S/cm $\approx$ 49 $\mu$S/cm. A conta da resolução está certa.%
\end{itemize}%
\paragraph{Pontos de dúvida}%
\begin{itemize}%
\item%
\textit{(solution)} A estimativa de ordem de grandeza usa p\_CO2 $\approx$ 1 atm. Mas o frasco já tinha ar a 1,0 atm, então a pressão parcial de CO2 no espaço gasoso é só \textasciitilde{}0,08 atm logo após a liberação e \textasciitilde{}0,01 atm no fim. Com esses valores, o mesmo modelo dá $\kappa$ de \textasciitilde{}14 $\mu$S/cm e \textasciitilde{}5 $\mu$S/cm, e não 50 $\mu$S/cm. Os \textasciitilde{}60 $\mu$S/cm do gráfico correspondem a uma situação perto de 1 atm de CO2. A resolução apresenta p ? 1 atm como limite superior e a conta serve só para ordem de grandeza, então a conclusão não muda. Ainda assim, a frase 'coerente com os $\approx$ 60 $\mu$S/cm' supõe implicitamente que o CO2 está a \textasciitilde{}1 atm. As constantes (k\_H, Ka1, $\lambda$) não vêm da prova, e a resolução declara isso.%
\item%
\textit{(statement)} Os gráficos são esquemáticos ('elaboração Comvest'). Valores lidos no gráfico, como 1,08 atm, 1,01 atm e 60 $\mu$S/cm, são aproximados e não precisam fechar quantitativamente com um modelo de equilíbrio simples. A questão é qualitativa: basta notar que a pressão sobe e depois cai e que a condutividade cresce de forma monotônica.%
\end{itemize}%
\section*{Questão 9}%
\begin{longtable}{@{}p{0.06\textwidth} p{0.33\textwidth} p{0.33\textwidth} p{0.1\textwidth} p{0.08\textwidth}@{}}\toprule Item & Resolução & Gabarito & Veredito & Certo \\ \midrule
Questão 9 (alternativa única) & (B) Au$^3$?; o nanocubo maciço vira uma estrutura oca e porosa (nanogaiola) de liga Au--Ag & B & igual & ambos \\
\bottomrule\end{longtable}%
\begin{itemize}%
\item%
\textbf{Questão 9 (alternativa única))} Conferi o enunciado na página 6 do PDF e ele bate com a transcrição: potenciais Pt$^2$?/Pt +1,20; Au$^3$?/Au +0,93; Ag?/Ag +0,80; Cu$^2$?/Cu +0,34; Zn$^2$?/Zn -0,76 V. Em (b), a figura mostra o cubo maciço virando um cubo vazado. Com Au$^3$?, E$^\circ$ = 0,93 - 0,80 = +0,13 V > 0, então a reação é espontânea; $\Delta$G$^\circ$ = -3$\cdot$96485$\cdot$0,13 $\approx$ -37,6 kJ/mol, compatível com os $\approx$ -38 kJ da resolução. Com Cu$^2$?, E$^\circ$ = 0,34 - 0,80 = -0,46 V < 0, então a reação não ocorre e (c) e (d) caem. A troca de 3 Ag por 1 Au e o fato de os dois metais serem CFC com parâmetros de rede quase iguais (Au 4,078 Å; Ag 4,086 Å) justificam que sobra uma casca de liga oca e porosa, o que elimina (a). Resolução e gabarito estão certos.%
\end{itemize}%
\paragraph{Pontos de dúvida}%
\begin{itemize}%
\item%
\textit{(statement)} O valor dado para Au$^3$?/Au (+0,93 V) é bem menor que o tabelado para o íon Au$^3$? livre ($\approx$ +1,50 V). Ele se aproxima mais do par AuCl$_{4}$$^{-}$/Au ($\approx$ +0,93 a +1,00 V), que é a espécie usada de fato no experimento com HAuCl$_{4}$. Isso não muda a resposta, porque nos dois casos o potencial é maior que +0,80 V.%
\item%
\textit{(solution)} A resolução descreve o mecanismo de forma simplificada: o Au se deposita na superfície e a Ag é consumida pelo interior. Na literatura (nanogaiolas do grupo de Xia), a prata dissolve primeiro por um poro numa das faces; depois, a desliga (dealloying) da casca Au--Ag é que abre os poros. Para a questão de múltipla escolha, o raciocínio basta e está conceitualmente correto.%
\end{itemize}%
\section*{Questão 10}%
\begin{longtable}{@{}p{0.06\textwidth} p{0.33\textwidth} p{0.33\textwidth} p{0.1\textwidth} p{0.08\textwidth}@{}}\toprule Item & Resolução & Gabarito & Veredito & Certo \\ \midrule
10 & D (desfavorecida com o aumento da temperatura, mas favorecida com o aumento da pressão) & D & igual & ambos \\
\bottomrule\end{longtable}%
\begin{itemize}%
\item%
\textbf{10)} Conferi o enunciado na página 6 do PDF (página 4 do caderno) e ele bate com a transcrição. O gráfico mostra K caindo sem parar quando T aumenta: fica um pouco acima de 10\^{}5 perto de 250 K e chega a cerca de 10\^{}{-}4 perto de 880 K, com a curva marcada como 300 atm. Pela equação de van 't Hoff e pelo princípio de Le Chatelier, isso indica $\Delta$H$^\circ$ < 0, e aquecer desfavorece o NH3. Para a pressão, $\Delta$n\_gás = 1 - (1/2 + 3/2) = -1. Com p\_i = x\_i$\cdot$P, a razão das frações molares vale K$\cdot$P, que aumenta com P quando T é fixa. Então aumentar a pressão favorece o NH3. Os dois efeitos levam à alternativa d), a mesma da resolução e do gabarito oficial.%
\end{itemize}%
\paragraph{Pontos de dúvida}%
\begin{itemize}%
\item%
\textit{(statement)} O enunciado diz "De acordo com a figura", mas a figura traz uma curva só, a 300 atm. Por isso o efeito da pressão não sai do gráfico: ele vem da estequiometria ($\Delta$n\_gás = -1) e do princípio de Le Chatelier. A resolução explica isso corretamente. Fica só a observação de que a pergunta mistura a leitura da figura com um conhecimento externo a ela.%
\item%
\textit{(solution)} A resolução diz que K não depende da pressão e trata os gases como ideais (p\_i = x\_i$\cdot$P). A 300 atm o comportamento real se afasta do ideal, e o rótulo "300 atm" no gráfico sugere que K pode ter sido calculada com correções de fugacidade. Nada disso muda a conclusão qualitativa, mas a hipótese de gás ideal não foi declarada.%
\item%
\textit{(solution)} Os valores de K (cerca de 10\^{}5 a 250 K e cerca de 10\^{}{-}4 a 880 K) foram lidos no gráfico, numa escala logarítmica. Erros de leitura de meia década não afetam a resposta, porque só importa que a curva é decrescente.%
\end{itemize}%
\section*{Questão 11}%
\begin{longtable}{@{}p{0.06\textwidth} p{0.33\textwidth} p{0.33\textwidth} p{0.1\textwidth} p{0.08\textwidth}@{}}\toprule Item & Resolução & Gabarito & Veredito & Certo \\ \midrule
Questão 11 (única) & (c) A e B & C & igual & ambos \\
\bottomrule\end{longtable}%
\begin{itemize}%
\item%
\textbf{Questão 11 (única))} Conferi o enunciado nas páginas 6--7 do PDF e ele bate com a transcrição: A é o ácido 2{-}sulfotereftálico, B é o ácido tereftálico e C é o 2{-}metilimidazol. Refazendo o raciocínio: o ligante A tem um grupo --SO3H que não participa da coordenação e fica livre no poro. Esse grupo é muito polar e ácido, faz ligação de hidrogênio e interação íon--dipolo com a água, então A deve ser a primeira opção. B só tem carboxilatos coordenados ao metal, o que dá afinidade intermediária. Em C (ZIF{-}8), os dois N se coordenam ao Zn$^2$? e o --CH3 é apolar, por isso o material é hidrofóbico e fica por último. A ordem A > B > C leva à alternativa (c), a mesma do gabarito. A justificativa química da resolução está correta.%
\end{itemize}%
\paragraph{Pontos de dúvida}%
\begin{itemize}%
\item%
\textit{(solution)} No passo 3, a resolução escreve 'A (--SO3H + 2 --COOH)' e 'B (2 --COOH)'. Isso descreve o ligante livre, mas na MOF os --COOH viram carboxilatos coordenados ao metal, como a própria resolução diz no passo 2. A redação é um pouco imprecisa, mas não muda a ordem nem a resposta.%
\item%
\textit{(solution)} A resolução não cita que, no UiO{-}66 (B), os grupos $\mu$3{-}OH dos clusters Zr6 também contribuem para a afinidade por água. Isso vai além do que a questão pede, que trata só dos ligantes, e não muda a ordem A > B > C.%
\end{itemize}%
\section*{Questão 12}%
\begin{longtable}{@{}p{0.06\textwidth} p{0.33\textwidth} p{0.33\textwidth} p{0.1\textwidth} p{0.08\textwidth}@{}}\toprule Item & Resolução & Gabarito & Veredito & Certo \\ \midrule
Questão 12 & C: {[}CH3OH{]} > 50 mg/dL, acidose metabólica (HCO3{-} < 15 mmol/L e pH arterial < 7,30) & C & igual & ambos \\
\bottomrule\end{longtable}%
\begin{itemize}%
\item%
\textbf{Questão 12)} A resposta da resolução é igual à do gabarito. Conferi a transcrição com o original (p. 5 do caderno, p. 7 do PDF) e ela está fiel. O raciocínio está certo. O metanol vira HCOOH, e o H+ liberado é neutralizado pelo HCO3{-}, o que faz o bicarbonato e o pH caírem: é uma acidose metabólica. O tratamento é indicado com o tóxico acima do limite de 50 mg/dL. Só a alternativa (c) junta as três condições de forma coerente. A (a) e a (d) falam em alcalose, e a (b) e a (d) põem o metanol abaixo do limite. Refiz as contas com python: 50 mg/dL = 0,50 g/L, que dividido por 32,04 g/mol dá $\approx$ 15,6 mmol/L. Com bicarbonato normal, 6,1 + log(24/(0,03$\cdot$40)) = 7,40. O pKa do HCOOH é $\approx$ 3,74. As três contas batem com a resolução.%
\end{itemize}%
\paragraph{Pontos de dúvida}%
\begin{itemize}%
\item%
\textit{(statement)} O enunciado usa "acidez metabólica". O termo clínico correto é "acidose metabólica". A resolução percebe isso e trata os dois como a mesma coisa, o que não muda a resposta.%
\item%
\textit{(solution)} O item 3 diz que \textasciitilde{}16 mmol/L de metanol geram ácido suficiente para consumir boa parte do bicarbonato (\textasciitilde{}24 mmol/L). É só um argumento de ordem de grandeza: não leva em conta a cinética do metabolismo, o volume de distribuição nem outros tampões. Não afeta a resposta, que depende do sentido das desigualdades e não dessa conta.%
\end{itemize}%
\section*{Questão 13}%
\begin{longtable}{@{}p{0.06\textwidth} p{0.33\textwidth} p{0.33\textwidth} p{0.1\textwidth} p{0.08\textwidth}@{}}\toprule Item & Resolução & Gabarito & Veredito & Certo \\ \midrule
Q13 & (D) Gases H2 e Cl2; Cl2(g) + H2O(l) ? HClO(aq) + H+(aq) + Cl{-}(aq) & D & igual & ambos \\
\bottomrule\end{longtable}%
\begin{itemize}%
\item%
\textbf{Q13)} Conferi a transcrição na página 7 do PDF (página 5 do caderno) e ela está fiel, inclusive o 'H2O(aq)' da alternativa (b), que está assim no próprio caderno. A química está correta. No cátodo a água se reduz a H2 e OH{-}, porque E$^\circ$(Na+/Na) = -2,71 V. No ânodo o Cl{-} se oxida a Cl2: a sobretensão do O2 é alta, como no processo cloro{-}álcali. Na água, o Cl2 se desproporciona em Cl2 + H2O ? HClO + H+ + Cl{-}, e essa equação fecha em Cl (2 = 2), H (2 = 2), O (1 = 1) e carga (0 = 0). As demais alternativas foram descartadas com razão. Em (a), H e carga não fecham (2 H contra 1 H; carga 0 contra -1). Em (b), há 2 H contra 3 H e carga -1 contra 0. Em (c), há 3 O contra 1 O e carga -1 contra +1. Além disso, (b) e (c) indicam O2 como gás produzido. A resolução e o gabarito concordam e estão certos.%
\end{itemize}%
\paragraph{Pontos de dúvida}%
\begin{itemize}%
\item%
\textit{(solution)} Para explicar por que o Cl{-} se oxida em vez da água, a resolução compara os potenciais padrão 1,23 V e 1,36 V. Em água de piscina, com pH perto de 7, o potencial do par O2/H2O cai para cerca de 0,82 V. Pela termodinâmica, então, o O2 seria ainda mais favorecido. A justificativa por sobretensão continua valendo, e a conclusão não muda.%
\item%
\textit{(solution)} Numa piscina não há membrana separando os eletrodos. Por isso o Cl2 encontra o OH{-} formado no cátodo e também gera ClO{-} (Cl2 + 2 OH{-} $\to$ ClO{-} + Cl{-} + H2O), e as quantidades de HClO e ClO{-} dependem do pH. O Cl2 também tende a se dissolver em vez de escapar como gás. A resolução não discute isso, mas o enunciado pede a equação de formação do HClO, e a alternativa (d) é a única coerente.%
\end{itemize}%
\section*{Questão 14}%
\begin{longtable}{@{}p{0.06\textwidth} p{0.33\textwidth} p{0.33\textwidth} p{0.1\textwidth} p{0.08\textwidth}@{}}\toprule Item & Resolução & Gabarito & Veredito & Certo \\ \midrule
14 & (B): a absorção é mais rápida que a metabolização; com 35 mg/dia (0,7 mg/kg), atingir o mínimo de 120 mg/kg $\times$ 50 kg = 6000 mg leva 6000/35 $\approx$ 171 dias, ou seja, cerca de 170 dias. & B & igual & ambos \\
\bottomrule\end{longtable}%
\begin{itemize}%
\item%
\textbf{14)} Conferi com a figura da p. 7 (folha 5 do caderno). A curva tracejada sobe em degraus de cerca de 35 mg a cada 24 h, chegando a $\approx$350 mg (eixo direito) no intervalo de 216--240 h, então a dose é 35 mg/dia, ou 0,7 mg/kg, dentro da faixa de 0,5--1,0 da bula. Refiz as contas: 120$\cdot$50 = 6000 mg e 6000/35 = 171,4 dias ($\approx$170 dias). O valor de 215 dias vem do limite superior: 7500/35 = 214,3 dias. No gráfico, a concentração sobe rápido logo depois de cada dose e cai devagar, e os picos e vales se acumulam, então a entrada no plasma é mais rápida que a eliminação. Assim ficam descartadas a duração da (a) e da (c) e o sentido da (c) e da (d). A resolução e o gabarito oficial concordam: (B).%
\end{itemize}%
\paragraph{Pontos de dúvida}%
\begin{itemize}%
\item%
\textit{(statement)} Comparar a 'velocidade de absorção' com a 'de metabolização' só faz sentido para as fases de subida, ou para o acúmulo nas primeiras doses. No estado estacionário (a partir de $\approx$120 h), a média da entrada é igual à da eliminação. A resolução usa a subida rápida e o acúmulo como justificativa, que é a interpretação pretendida pela banca. Não muda a alternativa.%
\item%
\textit{(solution)} A dose diária de 35 mg foi lida no gráfico, nos degraus da curva tracejada (eixo direito). Mesmo com um erro pequeno de leitura (por exemplo, 34--36 mg), o tempo mínimo fica entre 167 e 176 dias, mais perto de 170 do que de 215. A alternativa não muda.%
\item%
\textit{(statement)} 'No mínimo' se refere à dose desta paciente (35 mg/dia), e não à dose máxima da bula (50 mg/dia, que daria 120 dias). Nenhuma alternativa traz 120 dias, então a leitura não é ambígua para escolher a resposta.%
\end{itemize}%
\section*{Questão 15}%
\begin{longtable}{@{}p{0.06\textwidth} p{0.33\textwidth} p{0.33\textwidth} p{0.1\textwidth} p{0.08\textwidth}@{}}\toprule Item & Resolução & Gabarito & Veredito & Certo \\ \midrule
Q15 & (a) 95 pizzas de brigadeiro (e 110 de calabresa) & A & igual & ambos \\
\bottomrule\end{longtable}%
\begin{itemize}%
\item%
\textbf{Q15)} A transcrição confere com a página 8 do PDF original. O sistema c+b=205 e 45c+40b=8750 dá 45$\cdot$205-8750=475, e 475/5=95, logo b=95. Conferi com python3: 45$\cdot$110+40$\cdot$95=8750. As outras alternativas dão 8725 (b=100), 8700 (b=105) e 8675 (b=110). A resolução cita corretamente que a alternativa d) (110) é o número de pizzas de calabresa, uma possível armadilha.%
\end{itemize}%
\section*{Questão 16}%
\begin{longtable}{@{}p{0.06\textwidth} p{0.33\textwidth} p{0.33\textwidth} p{0.1\textwidth} p{0.08\textwidth}@{}}\toprule Item & Resolução & Gabarito & Veredito & Certo \\ \midrule
16 & (D) Vermelho venceu com 2 pontos & D & igual & ambos \\
\bottomrule\end{longtable}%
\begin{itemize}%
\item%
\textbf{16)} Conferi a transcrição com a página 8 do PDF e ela bate: os mesmos pontos, a mesma regra de pontuação e as mesmas alternativas. Refiz as contas com frações exatas. As distâncias ao quadrado são: azul 5, 13/4, 13 e 17; vermelho 29/100, 4, 2 e 8. A pedra mais próxima é vermelha (r$^2$ = 0,29), logo vermelho vence. A azul mais próxima tem d$^2$ = 13/4 = 3,25. As vermelhas com r$^2$ < 3,25 são só (1/2, 1/5) e (-1, 1), o que dá 2 pontos. Nenhuma pedra fica exatamente à distância d, então não há empate. A resolução e o gabarito estão certos.%
\end{itemize}%
\section*{Questão 17}%
\begin{longtable}{@{}p{0.06\textwidth} p{0.33\textwidth} p{0.33\textwidth} p{0.1\textwidth} p{0.08\textwidth}@{}}\toprule Item & Resolução & Gabarito & Veredito & Certo \\ \midrule
17 & (B): y = det A(x) = {-}x\^{}3 + 27x + 10, uma cúbica com coeficiente líder negativo, mínimo local em ({-}3, {-}44), máximo local em (3, 64), intercepto y = 10 e raízes {-}5 e (5 $\pm$ $\surd$33)/2 $\approx$ {-}0,37 e 5,37 & B & igual & ambos \\
\bottomrule\end{longtable}%
\begin{itemize}%
\item%
\textbf{17)} Refiz o determinante com uma função 3$\times$3 em Python. Ele coincide com {-}x\^{}3+27x+10 em x = {-}5, {-}3, 0, 1, 2,5, 3 e 7, e quatro pontos já bastam para fixar uma cúbica. Também conferi os valores críticos: y({-}3) = {-}44, y(3) = 64 e y(0) = 10. As raízes são {-}5, {-}0,372 e 5,372. Comparei com a figura da página 8 do PDF original: a alternativa (b) vem de cima à esquerda, tem o mínimo abaixo do eixo x à esquerda do eixo y, corta o eixo x um pouco à esquerda da origem e o eixo y acima dela, atinge o máximo acima do eixo x à direita e depois desce. Isso descreve exatamente essa função. As alternativas (c) e (d) são quárticas e (a) tem coeficiente líder positivo, então ficam excluídas. A resolução está correta e bem justificada, e a resposta coincide com o gabarito.%
\end{itemize}%
\section*{Questão 18}%
\begin{longtable}{@{}p{0.06\textwidth} p{0.33\textwidth} p{0.33\textwidth} p{0.1\textwidth} p{0.08\textwidth}@{}}\toprule Item & Resolução & Gabarito & Veredito & Certo \\ \midrule
Q18 & (C) 10 & C & igual & ambos \\
\bottomrule\end{longtable}%
\begin{itemize}%
\item%
\textbf{Q18)} Pelas relações de Girard, r1$\cdot$r2$\cdot$r3 = {-}8 e b = {-}(r1+r2+r3). Fiz uma busca exaustiva das ternas inteiras com produto {-}8 e a menor soma é {-}10, com raízes ({-}1,{-}1,{-}8). Então b\_max = 10 e p(x) = (x+1)\^{}2(x+8) = x\^{}3+10x\^{}2+17x+8. O enunciado do PDF (pág. 9, impressa 7) bate com a transcrição.%
\end{itemize}%
\paragraph{Pontos de dúvida}%
\begin{itemize}%
\item%
\textit{(statement)} O enunciado diz que todas as soluções de p(x)=0 são inteiras, e a resolução lê isso como 'as três raízes complexas são inteiras'. Se 'soluções' fossem só as reais, raízes complexas conjugadas seriam permitidas e b não teria limite. A resolução declara essa hipótese e a justifica. É a leitura pretendida e não muda a alternativa.%
\end{itemize}%
\section*{Questão 19}%
\begin{longtable}{@{}p{0.06\textwidth} p{0.33\textwidth} p{0.33\textwidth} p{0.1\textwidth} p{0.08\textwidth}@{}}\toprule Item & Resolução & Gabarito & Veredito & Certo \\ \midrule
Questão 19 & (D) 07h30 e 19h30 (t = 7,5 h e t = 19,5 h, onde A = 0,6 m) & D & igual & ambos \\
\bottomrule\end{longtable}%
\begin{itemize}%
\item%
\textbf{Questão 19)} Conferi com o enunciado original (p. 9 do PDF): é A(t) = (7/10)$\cdot$sen(($\pi$/12)(2t+3)) + 13/10, com 0 $\leq$ t $\leq$ 24. Como a amplitude é positiva, o mínimo ocorre quando sen $\theta$ = -1, ou seja, $\theta$ = 3$\pi$/2 + 2k$\pi$. Para t em {[}0, 24{]}, $\theta$ vai de $\pi$/4 a 17$\pi$/4, e só servem $\theta$ = 3$\pi$/2 (2t+3 = 18, t = 7,5 h) e $\theta$ = 7$\pi$/2 (2t+3 = 42, t = 19,5 h). Período = 2$\pi$/($\pi$/6) = 12 h e A\_min = 0,6 m. Também conferi A(7) = 0,7$\cdot$sen(17$\pi$/12) + 1,3 $\approx$ 0,624 m, maior que 0,6 m. A resolução está correta e coincide com o gabarito D.%
\end{itemize}%
\section*{Questão 20}%
\begin{longtable}{@{}p{0.06\textwidth} p{0.33\textwidth} p{0.33\textwidth} p{0.1\textwidth} p{0.08\textwidth}@{}}\toprule Item & Resolução & Gabarito & Veredito & Certo \\ \midrule
20 & (C) 3/2 & C & igual & ambos \\
\bottomrule\end{longtable}%
\begin{itemize}%
\item%
\textbf{20)} A transcrição bate com o original (pág. 9 do PDF, Questão 20). A área de B é a soma das áreas dos retângulos de base 1 e altura f(n) = 3\^{}({-}n), com n = 0, 1, 2, \ldots{}. Isso dá a série geométrica 1 + 1/3 + 1/9 + \ldots{} = 1/(1 - 1/3) = 3/2. Conferi em python3: as somas parciais convergem para 1,5. A área exata de A, 1/ln 3 $\approx$ 0,9102, fica abaixo de 3/2, como deve ser numa estimativa por excesso. A análise dos distratores também está certa: 9/10 = 0,9 é menor que a área de A; 3/9 = 1/3 é a área do retângulo n = 1; 40/27 é a soma parcial de n = 0 a 3.%
\end{itemize}%
\section*{Questão 21}%
\begin{longtable}{@{}p{0.06\textwidth} p{0.33\textwidth} p{0.33\textwidth} p{0.1\textwidth} p{0.08\textwidth}@{}}\toprule Item & Resolução & Gabarito & Veredito & Certo \\ \midrule
21 & (A) 1 & A & igual & ambos \\
\bottomrule\end{longtable}%
\begin{itemize}%
\item%
\textbf{21)} O enunciado confere com o original (pág. 9 do PDF, página 7 da prova). Com a PA a, a+r, a+2r, a+3r, temos det M = a(a+3r) - (a+r)(a+2r) = -2r$^2$ = -8, logo r = $\pm$2. Como a razão é negativa, r = -2. Pelo traço, 2a + 3r = -4, o que dá a = 1. Conferi numericamente: a PA fica 1, -1, -3, -5, com det = -8 e traço = -4. Com r = +2 sairia a = -5 (alternativa d), que é o distrator descartado pela razão negativa. A resolução e o gabarito dão a alternativa A.%
\end{itemize}%
\section*{Questão 22}%
\begin{longtable}{@{}p{0.06\textwidth} p{0.33\textwidth} p{0.33\textwidth} p{0.1\textwidth} p{0.08\textwidth}@{}}\toprule Item & Resolução & Gabarito & Veredito & Certo \\ \midrule
22 & (B) \{(x,y) $\in$ $\mathbb{R}$$^2$: |x| $\leq$ y $\leq$ 2 - |x|\} & B & igual & ambos \\
\bottomrule\end{longtable}%
\begin{itemize}%
\item%
\textbf{22)} Conferi o enunciado na página 9 do PDF (página 7 do caderno), e ele bate com a transcrição. O losango de vértices (0,0), (1,1), (0,2) e (-1,1) tem por borda de baixo o gráfico de y = |x| e por borda de cima o de y = 2 - |x|. Por isso a região é |x| $\leq$ y $\leq$ 2 - |x|, o que equivale a |x| + |y - 1| $\leq$ 1. Fiz também uma conferência numérica com 200 mil pontos aleatórios: a alternativa (b) coincide com o quadrado em todos os pontos testados. As alternativas (a), (c) e (d) divergem em milhares de pontos. Os contraexemplos da resolução estão certos: (a) aceita (2,1), porque 2 $\geq$ 1 $\geq$ 0; (c) é vazia, porque exigiria 0 $\geq$ 2; (d) aceita (5,6), porque 5 $\leq$ 6 $\leq$ 7.%
\end{itemize}%
\section*{Questão 23}%
\begin{longtable}{@{}p{0.06\textwidth} p{0.33\textwidth} p{0.33\textwidth} p{0.1\textwidth} p{0.08\textwidth}@{}}\toprule Item & Resolução & Gabarito & Veredito & Certo \\ \midrule
Q23 & (C) {[}1,2{]} & C & igual & ambos \\
\bottomrule\end{longtable}%
\begin{itemize}%
\item%
\textbf{Q23)} O enunciado confere com o original (pág. 9 do PDF, numerada 7). Como f(a) - f(b) = (a - b)(a + b - 1), um intervalo deixa f injetora se, e somente se, não tiver dois pontos distintos com a + b = 1, ou seja, se x = 1/2 não for ponto interior dele. Refiz as contas: f(0) = f(1) = -2 e f(-1) = f(2) = 0. Isso elimina a), b) e d). Em {[}1,2{]}, f é estritamente crescente, com imagem {[}-2,0{]} e inversa f$^{-}$$^{1}$(y) = 1/2 + $\surd$(y + 9/4). Essa inversa leva -2 em 1 e 0 em 2. A alternativa C está correta.%
\end{itemize}%
\paragraph{Pontos de dúvida}%
\begin{itemize}%
\item%
\textit{(statement)} O enunciado não diz qual é B. A resolução supõe, sem declarar, que B pode ser escolhido como f(A) = {[}-2,0{]}, o que garante a sobrejetividade. É a leitura natural da expressão "uma possibilidade para o conjunto A" e não muda a resposta.%
\end{itemize}%
\section*{Questão 24}%
\begin{longtable}{@{}p{0.06\textwidth} p{0.33\textwidth} p{0.33\textwidth} p{0.1\textwidth} p{0.08\textwidth}@{}}\toprule Item & Resolução & Gabarito & Veredito & Certo \\ \midrule
24 & (D) 149/198 & D & igual & ambos \\
\bottomrule\end{longtable}%
\begin{itemize}%
\item%
\textbf{24)} Conferi o enunciado na página 10 do PDF (Questão 24), e a transcrição está fiel. O mdc é par se e somente se os dois números são pares. Então P(mdc par) = C(50,2)/C(100,2) = 1225/4950 = 49/198, e P(mdc ímpar) = 1 - 49/198 = 149/198. Uma força bruta em Python, com mdc de todos os 4950 pares, dá 3725 casos favoráveis, ou seja, 149/198. A resolução e o gabarito batem.%
\end{itemize}%
\section*{Questão 25}%
\begin{longtable}{@{}p{0.06\textwidth} p{0.33\textwidth} p{0.33\textwidth} p{0.1\textwidth} p{0.08\textwidth}@{}}\toprule Item & Resolução & Gabarito & Veredito & Certo \\ \midrule
25 & (D) 1/16 < A < 3/32 & D & igual & ambos \\
\bottomrule\end{longtable}%
\begin{itemize}%
\item%
\textbf{25)} Confere com o enunciado original (p. 10 do PDF, Questão 25). Q e R têm a mesma abscissa, x = 1/2. Tomando QR como base, ela mede 1/2 - f(1/2), e a altura até P vale 1/2. Logo A = (1/4)(1/2 - f(1/2)). Pela hipótese com x = 1/2, 1/8 < f(1/2) < 1/4. Conferi pelo determinante, com frações exatas: nos extremos f(1/2) = 1/8 e f(1/2) = 1/4, a área vale 3/32 e 1/16. Como as desigualdades são estritas, 1/16 < A < 3/32, que é a alternativa D.%
\end{itemize}%
\section*{Questão 26}%
\begin{longtable}{@{}p{0.06\textwidth} p{0.33\textwidth} p{0.33\textwidth} p{0.1\textwidth} p{0.08\textwidth}@{}}\toprule Item & Resolução & Gabarito & Veredito & Certo \\ \midrule
26 & (a) AP = $\surd$(2 - $\surd$2) $\approx$ 0,765 & A & igual & ambos \\
\bottomrule\end{longtable}%
\begin{itemize}%
\item%
\textbf{26)} O enunciado confere com o original (pág. 10 do PDF). Refiz a conta. O triângulo OCB é isósceles (OC = CB = 1), então $\angle$COB = 15$^\circ$ e $\angle$OCB = 150$^\circ$. Daí $\angle$OCA = 30$^\circ$, e como OAC também é isósceles, $\angle$OAC = 30$^\circ$. Como P, O, B são colineares, $\angle$AOP é ângulo externo do triângulo AOB: $\angle$AOP = 30$^\circ$ + 15$^\circ$ = 45$^\circ$. Assim, AP = 2$\cdot$sen(22,5$^\circ$) = $\surd$(2 - $\surd$2) $\approx$ 0,7654. Também conferi numericamente: com B = (2cos15$^\circ$, 0), a reta de inclinação 15$^\circ$ corta a circunferência em C (ângulo de 15$^\circ$, com BC = 1) e em A (ângulo de 135$^\circ$), o que dá AP $\approx$ 0,7654. A resolução e o gabarito dão a mesma alternativa (a), e ela está correta.%
\end{itemize}%
\paragraph{Pontos de dúvida}%
\begin{itemize}%
\item%
\textit{(solution)} Na análise das alternativas, a resolução erra dois valores aproximados. $\surd$(4 - $\surd$2) dá $\approx$ 1,608, e não 1,612. $\surd$(4 + $\surd$2) dá $\approx$ 2,327, e não 2,298. Isso não muda nada: a alternativa (d) continua maior que o diâmetro 2, e a resposta (a) se mantém.%
\end{itemize}%
\section*{Questão 27}%
\begin{longtable}{@{}p{0.06\textwidth} p{0.33\textwidth} p{0.33\textwidth} p{0.1\textwidth} p{0.08\textwidth}@{}}\toprule Item & Resolução & Gabarito & Veredito & Certo \\ \midrule
27 & (D) trajado, em termos figurativos, algo que se apresenta como um imperativo moral & D & igual & ambos \\
\bottomrule\end{longtable}%
\begin{itemize}%
\item%
\textbf{27)} Conferi com o PDF original (p. 10, numeração impressa 8): a transcrição está fiel. O título do texto 2, "'FANTASIA' NOTA:10", usa a palavra entre aspas e vem do julgamento dos desfiles de carnaval, em que o quesito "fantasia" é o traje. Isso remete à acepção 4 do Houaiss (traje), usada em sentido figurado: a "roupa" do folião é a frase escrita no corpo, uma ordem no imperativo ("Não mate") com valor de dever moral. As alternativas a, b e c correspondem às acepções 1, 2 e 3 (imaginação, obra imaginária, ficção). Elas não servem, porque a mensagem trata de uma violência real e de algo "óbvio", como diz a legenda do texto 1. A resolução justifica bem cada alternativa e chega à mesma letra do gabarito.%
\end{itemize}%
\section*{Questão 28}%
\begin{longtable}{@{}p{0.06\textwidth} p{0.33\textwidth} p{0.33\textwidth} p{0.1\textwidth} p{0.08\textwidth}@{}}\toprule Item & Resolução & Gabarito & Veredito & Certo \\ \midrule
28 & (B) ao retorno à realidade com o término da festa & B & igual & ambos \\
\bottomrule\end{longtable}%
\begin{itemize}%
\item%
\textbf{28)} Conferi o enunciado nas páginas 10 e 11 do PDF, e a transcrição bate com a letra da canção e com as alternativas. A leitura da resolução está correta. "Desengano" é a perda de uma ilusão. O refrão "Quarta{-}feira sempre desce o pano" é uma metáfora teatral: a cortina desce, o carnaval acaba e volta a vida comum. Com isso volta também a dor que o eu lírico deixou "em casa me esperando" e se desfaz o romance passageiro ("hoje nem lembra não"). As outras alternativas não servem. A confraternização da (a) aparece ligada à "esperança", não ao desengano. O amor da (c) é breve e já esquecido, e o desengano vem do fim dele, não de tê{-}lo vivido. Na (d), a dor fica suspensa durante a festa, não aumentada.%
\end{itemize}%
\section*{Questão 29}%
\begin{longtable}{@{}p{0.06\textwidth} p{0.33\textwidth} p{0.33\textwidth} p{0.1\textwidth} p{0.08\textwidth}@{}}\toprule Item & Resolução & Gabarito & Veredito & Certo \\ \midrule
29 & (D) presente do subjuntivo (viva, possa e saiba, nos versos 12, 13 e 14) & D & igual & ambos \\
\bottomrule\end{longtable}%
\begin{itemize}%
\item%
\textbf{29)} Conferi a transcrição com o original (págs. 10 e 11 do PDF) e ela está fiel. O texto de Módolo e Braga liga o futuro à expectativa, e a questão pede formas verbais que exprimam expectativa pelo que "ainda não é". Na última estrofe, aberta por "No carnaval, esperança", as orações optativas "Que gente longe viva\ldots{} / possa entrar\ldots{} / saiba ser\ldots{}" trazem verbos no presente do subjuntivo, modo do desejo e do que ainda não se realizou. As outras alternativas não servem: (a) pretérito perfeito, ligado à memória; (b) gerúndios ligados a ações passadas (deixei, deixou); (c) presente do indicativo habitual ou de constatação, ligado à atenção. A resolução justifica bem a escolha e descarta corretamente cada distrator.%
\end{itemize}%
\paragraph{Pontos de dúvida}%
\begin{itemize}%
\item%
\textit{(solution)} Na alternativa (b), "a dor em casa me esperando" tem sentido lexical de espera e de antecipação do que virá (a volta à dor depois do carnaval). A resolução descarta a alternativa porque o gerúndio acompanha fatos passados e não marca projeção. O argumento convence, mas um professor pode lembrar que o distrator foi feito para atrair pelo significado da palavra, e a questão pede o tempo/modo verbal, não o sentido lexical.%
\item%
\textit{(solution)} Em (c), "Quarta{-}feira sempre desce o pano" é um presente habitual (gnômico) que também aponta para um fim previsível. A resolução o classifica como constatação ligada à atenção, o que é razoável e não muda a resposta.%
\end{itemize}%
\section*{Questão 30}%
\begin{longtable}{@{}p{0.06\textwidth} p{0.33\textwidth} p{0.33\textwidth} p{0.1\textwidth} p{0.08\textwidth}@{}}\toprule Item & Resolução & Gabarito & Veredito & Certo \\ \midrule
30 & (A) postura discriminatória disfarçada de enaltecimento & A & igual & ambos \\
\bottomrule\end{longtable}%
\begin{itemize}%
\item%
\textbf{30)} Reli o enunciado no PDF (p. 11 do arquivo, página 9 do caderno), e a transcrição está fiel. O próprio texto explica o termo: "Você passa bem café" é um louvor{-}prisão, "um rebaixamento naquele universo onde o valor era dado pelos livros". O termo é dito "elogio sutil com uma etiqueta de gênero" e "mais agressivo" que o insulto direto. No fim do trecho, o autor o resume como "agredir com um elogio, preconceito com açúcar, ataque com punho de renda". Isso é discriminação disfarçada de elogio, ou seja, a alternativa A. A alternativa B corresponde ao "veado", o ataque direto e explícito, que o texto contrapõe ao louvor{-}prisão. A C não serve: o texto não fala de empatia e diz que o elogio é mais agressivo. A D descreve a reação de quem recebe, e não o elogio. Além disso, o texto liga a resposta defensiva ao primeiro insulto. A resolução descarta as três corretamente e chega à mesma letra do gabarito.%
\end{itemize}%
\section*{Questão 31}%
\begin{longtable}{@{}p{0.06\textwidth} p{0.33\textwidth} p{0.33\textwidth} p{0.1\textwidth} p{0.08\textwidth}@{}}\toprule Item & Resolução & Gabarito & Veredito & Certo \\ \midrule
31 & C --- os dois textos recorrem a informações estatísticas na defesa de um ponto de vista & C & igual & ambos \\
\bottomrule\end{longtable}%
\begin{itemize}%
\item%
\textbf{31)} A transcrição confere com o original (págs. 11--12 do PDF). O Texto 1 sustenta que o brasileiro trabalha menos do que o esperado e se apoia na pesquisa com 160 países e nas médias de 42,7 h e 40,1 h por semana. O Texto 2 sustenta a tese contrária com o resíduo econométrico, os controles institucionais e os microdados da World Values Survey (58 países). As teses são opostas, mas os dois usam estatística para defender o que dizem, o que confirma a alternativa C. A resolução também descarta bem as outras. Na (b), só o Texto 2 contesta o argumento da mídia, e o Texto 1 é justamente a matéria que o divulga. Na (d), a divergência existe entre os dois textos, mas nenhum deles, isoladamente, apresenta uma divergência entre especialistas. Na (a), nenhum dos textos critica aspectos culturais e econômicos: o Texto 2 critica a metodologia.%
\end{itemize}%
\paragraph{Pontos de dúvida}%
\begin{itemize}%
\item%
\textit{(solution)} A resolução escreve o nome do autor do Texto 1 como "Carriello", mas a fonte da prova traz "CARRIELO". É só um erro de digitação, sem efeito na resposta.%
\end{itemize}%
\section*{Questão 32}%
\begin{longtable}{@{}p{0.06\textwidth} p{0.33\textwidth} p{0.33\textwidth} p{0.1\textwidth} p{0.08\textwidth}@{}}\toprule Item & Resolução & Gabarito & Veredito & Certo \\ \midrule
32 & B (reverter a conclusão à qual chega o texto 1) & B & igual & ambos \\
\bottomrule\end{longtable}%
\begin{itemize}%
\item%
\textbf{32)} A transcrição bate com o original (págs. 9--10 do caderno). Pelo texto 1, mesmo considerando demografia, renda, impostos e transferências, o brasileiro trabalha menos do que seria esperado, e o motivo provável seria uma preferência cultural pelo lazer. O texto 2 usa os mesmos dados de Gethin e Saez. Com controles institucionais, mostra que 'o Brasil passa a trabalhar mais do que o esperado' e rejeita a hipótese cultural. Por isso a conclusão do texto 1 é revertida (B). A alternativa (a) não serve, porque o texto 2 contesta os resultados em vez de completá{-}los. A (c) também não, porque o banco de dados não é invalidado: o próprio texto 2 o usa, e a crítica é à especificação econométrica. A (d) cai porque a hipótese é rejeitada, não confirmada. A análise de cada alternativa na resolução está correta e fiel ao texto.%
\end{itemize}%
\section*{Questão 33}%
\begin{longtable}{@{}p{0.06\textwidth} p{0.33\textwidth} p{0.33\textwidth} p{0.1\textwidth} p{0.08\textwidth}@{}}\toprule Item & Resolução & Gabarito & Veredito & Certo \\ \midrule
Q33 & (C) caracterização das relações hierárquicas, pois o poder senhorial se exerce sobre os subalternizados & C & igual & ambos \\
\bottomrule\end{longtable}%
\begin{itemize}%
\item%
\textbf{Q33)} Conferi o enunciado na página 12 do PDF (página 10 do caderno), e ele bate com a transcrição. A resolução e o gabarito concordam e a interpretação está certa. O trecho expõe a relação senhor--escravizado sem disfarce: Prudêncio é animalizado ("cavalo", "freio", "besta"), apanha ("fustigava{-}o", "gemendo") e é silenciado ("obedecia sem dizer palavra", "Cala a boca, besta!"). Isso caracteriza o poder senhorial exercido sobre o subalternizado. As outras alternativas não se sustentam. A (a) toma a "brincadeira" como atenuação, quando é justamente a forma em que a dominação se exerce. Na (b), o poder nas mãos de uma criança não enfraquece o quadro social; mostra que a dominação já está naturalizada. Na (d), o "ai, nhonhô!" é queixa logo reprimida, e não inversão de papéis. A resolução cita corretamente o capítulo "O vergalho": Prudêncio, já liberto, chicoteia um escravizado que comprou e repete "Cala a boca, besta!". Isso reforça que a hierarquia se reproduz, e não se inverte.%
\end{itemize}%
\paragraph{Pontos de dúvida}%
\begin{itemize}%
\item%
\textit{(solution)} Erro de digitação em LaTeX: "se \textbackslash{}emph\{reproduz*;" deveria ser "\textbackslash{}emph\{reproduz\}". A chave sem fechamento pode quebrar a compilação ou deixar o resto do parágrafo em itálico. O conteúdo não muda.%
\item%
\textit{(solution)} Na análise da (b), a resolução diz que "o narrador não denuncia nada: ele relata com indiferença". No passo 2, porém, atribui à "ironia machadiana" a função de expor a brutalidade. Não há contradição: um é o narrador, o outro é o autor implícito. Ainda assim, a redação poderia separar melhor narrador e autor.%
\end{itemize}%
\section*{Questão 34}%
\begin{longtable}{@{}p{0.06\textwidth} p{0.33\textwidth} p{0.33\textwidth} p{0.1\textwidth} p{0.08\textwidth}@{}}\toprule Item & Resolução & Gabarito & Veredito & Certo \\ \midrule
34 & D & D & igual & ambos \\
\bottomrule\end{longtable}%
\begin{itemize}%
\item%
\textbf{34)} A transcrição bate com o original (p. 12 do PDF). A resposta D está certa. A 'advertência' é um recurso ficcional, o do editor ou autor fictício: Lima Barreto se apresenta como simples editor e atribui a autoria a Augusto Machado, que é o narrador{-}personagem do romance. A resolução descarta bem as outras alternativas. A (a) inverte a relação, porque o autor real é Lima Barreto, e o trecho faz parte da ficção. Na (b), o 'editor' não critica os defeitos, chama{-}os de insignificantes e de desprezíveis, e o próprio trecho diz que o autor{-}narrador 'sobressai demasiado', o que contradiz a ideia de 'discrição e objetividade'. Na (c), o Modernismo não separou os gêneros, e a obra mistura biografia e romance. Os dados de apoio também estão corretos: publicação em 1919 e o mesmo recurso usado por Machado de Assis com o Conselheiro Aires em Esaú e Jacó e no Memorial de Aires.%
\end{itemize}%
\section*{Questão 35}%
\begin{longtable}{@{}p{0.06\textwidth} p{0.33\textwidth} p{0.33\textwidth} p{0.1\textwidth} p{0.08\textwidth}@{}}\toprule Item & Resolução & Gabarito & Veredito & Certo \\ \midrule
35 & (C) expressam o protagonismo daquele que conduz, condição recuperada pela aproximação semântica entre as palavras bengala e cajado. & C & igual & ambos \\
\bottomrule\end{longtable}%
\begin{itemize}%
\item%
\textbf{35)} A transcrição confere com o original (p. 10 do caderno, página 12 do PDF). A resposta da resolução coincide com o gabarito oficial. A leitura está certa. Os passos do eu{-}lírico formam o rebanho e ele vira o pastor. A bengala ocupa o lugar do cajado, o bastão do pastor, por aproximação semântica. Assim, o que costuma lembrar velhice e dependência passa a indicar quem conduz ('me faço pastor'). As outras alternativas foram bem descartadas. Em (a), não há 'dificuldade de comparação' nem 'fantasia'. Em (b), o texto não fala de trauma, e não há 'passos do rebanho': os passos do eu{-}lírico é que são o rebanho. Em (d), não há lamento nem 'dissolução sintática', já que o poema é um único período completo e a troca bengala/cajado acontece no sentido, não na sintaxe.%
\end{itemize}%
\paragraph{Pontos de dúvida}%
\begin{itemize}%
\item%
\textit{(solution)} A resolução diz que 'o livro chama{-}se Odes Mínimas'. O livro citado é 'Prosas seguidas de Odes Mínimas', e 'Odes Mínimas' é só a seção onde está o poema. A imprecisão é pequena e não afeta a resposta.%
\item%
\textit{(solution)} A resolução trata a ode como gênero de 'exaltação' e usa isso para descartar o lamento na alternativa (d). As 'Odes Mínimas' de José Paulo Paes celebram objetos do dia a dia, muitas vezes com humor ou ironia. O argumento serve, mas a conclusão já se sustenta só pelo texto: com a bengala, o eu{-}lírico 'se faz' pastor, o que é uma postura ativa e não de lamento.%
\end{itemize}%
\section*{Questão 36}%
\begin{longtable}{@{}p{0.06\textwidth} p{0.33\textwidth} p{0.33\textwidth} p{0.1\textwidth} p{0.08\textwidth}@{}}\toprule Item & Resolução & Gabarito & Veredito & Certo \\ \midrule
Questão 36 (alternativa única) & D & D & igual & ambos \\
\bottomrule\end{longtable}%
\begin{itemize}%
\item%
\textbf{Questão 36 (alternativa única))} A transcrição confere com o original (p. 11 do caderno, p. 13 do PDF). No texto 1, a dor da autora está no ato de escrever ("escrever 'Zaíta' é um processo doloroso"), não no que os contos narram. Isso descarta (c). As alternativas (a) e (b) dizem que a escrita "diminui" o abandono ou o silêncio familiar, o que não faz sentido. Além disso, os silêncios continuam presentes nos dois excertos: a mãe que "decifrava os seus silêncios" e Benícia, que "recomendou então o silêncio". O que os textos têm em comum é a infância e a maternidade sob pressão da miséria e da violência, que viram matéria literária pela escrita comprometida e dolorosa da autora. Isso corresponde a (d), a mesma resposta do gabarito oficial. A justificativa da resolução está correta e bem fundamentada nos excertos.%
\end{itemize}%
\paragraph{Pontos de dúvida}%
\begin{itemize}%
\item%
\textit{(solution)} A resolução usa o conceito de "escrevivência" e detalhes da leitura integral dos contos que não estão nos excertos: a fome disfarçada com brincadeiras, o filho envolvido com o tráfico e a morte de Zaíta por bala perdida num tiroteio. Esses detalhes estão de acordo com os contos e com o que o enunciado pede ("leitura prévia e integral"), mas não dá para conferi{-}los só com o caderno. Não mudam a alternativa escolhida.%
\item%
\textit{(solution)} Para descartar (a), a resolução diz que não há abandono familiar porque "a mãe está presente". O argumento principal, de que a escrita não "diminui" nada dentro das narrativas, já basta. A afirmação sobre ausência de abandono é discutível, porque a figura paterna é ausente nos contos, mas isso não altera a conclusão.%
\end{itemize}%
\section*{Questão 37}%
\begin{longtable}{@{}p{0.06\textwidth} p{0.33\textwidth} p{0.33\textwidth} p{0.1\textwidth} p{0.08\textwidth}@{}}\toprule Item & Resolução & Gabarito & Veredito & Certo \\ \midrule
Q37 & (B) Ao mudar o nome para Afamefuna, Grace se reconecta com as expectativas e os valores da avó, ligados à tradição anterior à colonização. & B & igual & ambos \\
\bottomrule\end{longtable}%
\begin{itemize}%
\item%
\textbf{Q37)} Conferi o enunciado na página 13 do PDF (página 11 do caderno) e a transcrição está fiel. A resolução identifica bem o contexto do conto. A avó Nwamgba deu à neta o nome igbo Afamefuna ("meu nome não se perderá"). O pai, convertido, passou a se chamar Michael, e a menina foi criada como Grace. No fim do excerto, Grace se sente "distante de suas raízes" e troca oficialmente Grace por Afamefuna, o que confirma (b). A resolução descarta as outras com argumentos sólidos. Em (a), o protagonismo de Grace vai contra a ocidentalização, não em paralelo a ela. Em (c), o sentido se inverte: o texto não trata o passado africano como "atraso" nem exalta a educação ocidental. Em (d), a anáfora "Foi Grace quem" resume a trajetória da personagem, não antecipa a tristeza. Além disso, o conflito é com a educação colonial, que nega a história africana, e não entre "avanço da educação" e "ignorância".%
\end{itemize}%
\paragraph{Pontos de dúvida}%
\begin{itemize}%
\item%
\textit{(solution)} A resolução usa detalhes do conto que não estão no excerto: Nwamgba, Michael, o significado de Afamefuna e o nome ter sido dado pela avó. Esses detalhes estão corretos e a própria alternativa (b) os pressupõe, mas não podem ser verificados só pelo trecho da prova. A frase "a educação ocidental aparece como algo que pode apagar a identidade africana" é uma leitura um pouco forte, porque a própria Grace se formou numa universidade de molde ocidental. Isso não muda a resposta.%
\end{itemize}%
\section*{Questão 38}%
\begin{longtable}{@{}p{0.06\textwidth} p{0.33\textwidth} p{0.33\textwidth} p{0.1\textwidth} p{0.08\textwidth}@{}}\toprule Item & Resolução & Gabarito & Veredito & Certo \\ \midrule
38 & C & C & igual & ambos \\
\bottomrule\end{longtable}%
\begin{itemize}%
\item%
\textbf{38)} Conferi o enunciado na p. 13 do PDF (página 11 da prova) e a transcrição está fiel. No conto, José Montiel enriquece comprando a preço vil os bens dos expulsos. A esposa, alheia a isso, acha que ele está ajudando os perseguidos e que vai se arruinar por caridade. O narrador mostra os dois lados: primeiro revela o negócio real ("Na verdade...") e depois põe a fala ingênua da mulher. Disso nasce uma ironia dramática: o leitor sabe o que a personagem ignora, que é exatamente a alternativa (c). Os descartes também estão certos. Em (a), as oposições são antíteses, não paradoxo. Em (b), o raciocínio da esposa é coerente, só parte de uma premissa falsa. Em (d), ela fala com sinceridade e não sabe das intenções do marido, então não há ironia dela nem efeito cômico.%
\end{itemize}%
\paragraph{Pontos de dúvida}%
\begin{itemize}%
\item%
\textit{(solution)} A resolução chama o recurso de "ironia dramática ou situacional". O termo exato para a situação descrita em (c), em que o leitor sabe o que a personagem ignora, é ironia dramática. A ironia situacional (quando o resultado contraria o esperado) é outro conceito. A imprecisão é só de nomenclatura e não muda a alternativa.%
\end{itemize}%
\section*{Questão 39}%
\begin{longtable}{@{}p{0.06\textwidth} p{0.33\textwidth} p{0.33\textwidth} p{0.1\textwidth} p{0.08\textwidth}@{}}\toprule Item & Resolução & Gabarito & Veredito & Certo \\ \midrule
Q39 & (B) por disputas de classificações, abarcadas nas diferentes ideologias raciais brasileiras, com impacto na produção de políticas públicas e na forma como os indivíduos se percebem. & B & igual & ambos \\
\bottomrule\end{longtable}%
\begin{itemize}%
\item%
\textbf{Q39)} A resolução e o gabarito marcam a mesma alternativa. Conferi o enunciado na página 14 do PDF e ele bate com a transcrição. O Texto 2 afasta explicitamente três explicações: o crescimento vegetativo diferencial, as mudanças na pergunta do censo e a hetero ou autoclassificação. Isso elimina (d). O texto também dá mais peso às ideologias raciais (embranquecimento, mestiçagem, negritude) do que aos incentivos materiais das políticas públicas, o que elimina (c). Por fim, diz que essas ideologias convivem sem que uma suplante as outras, o que contradiz a ideia de ideologia 'hegemônica' produzida pelo Estado da alternativa (a). Sobra (b). As leituras do gráfico na resolução estão corretas e não interferem na escolha: brancos $\approx$38\%$\to$63\%$\to$43\%, pardos $\approx$38\%$\to$21\%$\to$45\% e pretos $\approx$20\%$\to$15\%$\to$5--6\%$\to$10\%.%
\end{itemize}%
\paragraph{Pontos de dúvida}%
\begin{itemize}%
\item%
\textit{(solution)} A resolução diz que o período de 1872--1940 foi o 'auge da ideologia do embranquecimento' e liga a mestiçagem e a negritude a fases específicas do gráfico. Essas ligações vêm de conhecimento externo e não estão nos textos. São razoáveis, mas a resposta não depende delas.%
\item%
\textit{(statement)} A alternativa (b) usa as palavras 'disputas de classificações' e 'impacto na produção de políticas públicas'. O Texto 2 só diz que as ideologias 'conviveram e ainda convivem entre si, incapazes de suplantar umas às outras' e não trata do impacto sobre as políticas públicas. Para chegar a (b), é preciso inferir um pouco além do texto, como a resolução faz. Mesmo assim, as outras alternativas contradizem o texto, e (b) continua sendo a única defensável.%
\end{itemize}%
\section*{Questão 40}%
\begin{longtable}{@{}p{0.06\textwidth} p{0.33\textwidth} p{0.33\textwidth} p{0.1\textwidth} p{0.08\textwidth}@{}}\toprule Item & Resolução & Gabarito & Veredito & Certo \\ \midrule
Q40 & D & D & igual & ambos \\
\bottomrule\end{longtable}%
\begin{itemize}%
\item%
\textbf{Q40)} Conferi o enunciado na página 12 do caderno (página 14 do PDF) e ele bate com a transcrição. A tese de Przeworski é que, sob polarização, a democracia 'só pode ser defendida como o melhor método para administrar essas divergências' e que seu 'milagre' é administrar os conflitos 'sem repressão e em paz'. Isso corresponde à alternativa D. As alternativas A e B introduzem a ideia de 'maioria' e de crise ou questionamento, que não aparecem no texto. A alternativa C contradiz o texto, que afirma justamente que os conflitos podem ser administrados sem repressão. A análise da resolução está correta.%
\end{itemize}%
\paragraph{Pontos de dúvida}%
\begin{itemize}%
\item%
\textit{(solution)} A alternativa D fala em 'método eficaz', mas o texto diz que a democracia 'só pode ser defendida como o melhor método' (é uma defesa normativa, não uma constatação de eficácia). A resolução trata as duas ideias como equivalentes sem comentar a diferença. Isso não muda a resposta, porque D é claramente a única alternativa compatível com o texto.%
\end{itemize}%
\section*{Questão 41}%
\begin{longtable}{@{}p{0.06\textwidth} p{0.33\textwidth} p{0.33\textwidth} p{0.1\textwidth} p{0.08\textwidth}@{}}\toprule Item & Resolução & Gabarito & Veredito & Certo \\ \midrule
Q41 & (C) se apropriam de retóricas e credenciais científicas para produzir falsas controvérsias que visam gerar dissensos e incertezas na sociedade & C & igual & ambos \\
\bottomrule\end{longtable}%
\begin{itemize}%
\item%
\textbf{Q41)} Conferi o enunciado na página 14 do PDF (impressa como 12), e a transcrição está fiel. A alternativa (c) repete quase palavra por palavra o trecho final do Texto 1: "apropriando{-}se deliberadamente de retóricas e credenciais científicas, com vistas a persuadir terceiros, criando dissensos e incertezas". Ela também se encaixa na ideia central do texto, a "fabricação de falsas controvérsias científicas". As demais alternativas foram bem descartadas. (a) contraria o trecho "não se trata de um mero repúdio 'pseudo' ou 'anticientífico'". (b) fala em oposição externa às ciências, quando o texto descreve uma região de sobreposição nas franjas das ciências. (d) inverte a relação, porque os dissensos são criados pelos negacionistas e não existiam antes. Além disso, ninguém nega a existência das ciências. A resolução e o gabarito batem e estão corretos.%
\end{itemize}%
\paragraph{Pontos de dúvida}%
\begin{itemize}%
\item%
\textit{(solution)} A resolução lê a figura do Texto 2 de forma interpretativa. Ela chama as "controvérsias" de legítimas e próprias das ciências e diz que a fronteira tracejada e inclinada indica uma separação pouco nítida. A leitura é razoável e combina com o Texto 1, mas a figura não traz esses sentidos de forma explícita. Isso não altera a resposta.%
\end{itemize}%
\section*{Questão 42}%
\begin{longtable}{@{}p{0.06\textwidth} p{0.33\textwidth} p{0.33\textwidth} p{0.1\textwidth} p{0.08\textwidth}@{}}\toprule Item & Resolução & Gabarito & Veredito & Certo \\ \midrule
Q42 & (A) transcrição, que ocorre no núcleo; o DNA serve de molde para um RNAm com bases complementares, que leva a informação ao citoplasma & A & igual & ambos \\
\bottomrule\end{longtable}%
\begin{itemize}%
\item%
\textbf{Q42)} A transcrição confere com o PDF original (página 15, impressa como 13). A tirinha mostra o DNA ditando uma mensagem. A 'secretária', uma fita simples de RNA, escreve a fita complementar e usa U no lugar de T. O pareamento confere: CACA$\to$GUGU, CTCT$\to$GAGA, ACT$\to$UGA (C$\to$G, A$\to$U, T$\to$A). Isso é transcrição, que nos eucariotos ocorre no núcleo, e o RNAm vai para o citoplasma, onde ficam os ribossomos. A alternativa (a) é a única coerente. A resolução descarta corretamente as demais: (b) e (d) chamam o processo de tradução e falam em síntese de 'fita de RNAt'; (c) põe a transcrição no ribossomo, com RNAt como molde. A resolução e o gabarito concordam e estão certos.%
\end{itemize}%
\paragraph{Pontos de dúvida}%
\begin{itemize}%
\item%
\textit{(solution)} A resolução simplifica alguns pontos, o que é razoável para o nível da prova. Ela não diz que a transcrição usa apenas uma das fitas do DNA (a fita{-}molde). Também não comenta que, nos eucariotos, o transcrito primário (pré{-}RNAm) passa por processamento no núcleo antes de ir ao citoplasma. Nada disso muda a resposta.%
\end{itemize}%
\section*{Questão 43}%
\begin{longtable}{@{}p{0.06\textwidth} p{0.33\textwidth} p{0.33\textwidth} p{0.1\textwidth} p{0.08\textwidth}@{}}\toprule Item & Resolução & Gabarito & Veredito & Certo \\ \midrule
43 & (C) & C & igual & ambos \\
\bottomrule\end{longtable}%
\begin{itemize}%
\item%
\textbf{43)} Conferi a transcrição com a página 13 do caderno (página 15 do PDF) e ela está fiel. O texto diz que as toxinas sofreram seleção positiva. Isso quer dizer que a variante vantajosa (a duplicação gênica, que aumenta a letalidade para as aves) aumentou a sobrevivência e/ou a reprodução das serpentes que a carregam e ficou mais frequente ao longo do tempo, que é exatamente o que diz a alternativa (c). As outras estão erradas. (a) Fluxo gênico é a troca de alelos entre populações. Não é causado pela seleção natural e, numa população isolada na ilha, acontece pouco ou nada. (b) A diferenciação levou cerca de 100 mil anos, então não foi súbita, e uma duplicação isolada não basta para explicar a especiação. (d) A deriva genética muda as frequências alélicas ao acaso, mas não causa mutações como a duplicação gênica. A resolução identifica e justifica corretamente os erros de cada alternativa.%
\end{itemize}%
\paragraph{Pontos de dúvida}%
\begin{itemize}%
\item%
\textit{(solution)} Na alternativa (b), a resolução diz que 'mutações não são motivadas por uma necessidade'. Só que a alternativa diz outra coisa: que a duplicação 'motivou a especiação', ou seja, que a mutação foi a causa da especiação, e não algo motivado. A resolução leu esse trecho de forma um pouco equivocada. Mesmo assim, o descarte da (b) continua correto pelo outro argumento que ela dá: a especiação não foi súbita (foi ao longo de cerca de 100 mil anos) e envolveu isolamento e seleção. A resposta final não muda.%
\end{itemize}%
\section*{Questão 44}%
\begin{longtable}{@{}p{0.06\textwidth} p{0.33\textwidth} p{0.33\textwidth} p{0.1\textwidth} p{0.08\textwidth}@{}}\toprule Item & Resolução & Gabarito & Veredito & Certo \\ \midrule
44 & B: (i) cnidários; (ii) platelmintos; (iii) da membrana plasmática; (iv) sódio; (v) potássio & B & igual & ambos \\
\bottomrule\end{longtable}%
\begin{itemize}%
\item%
\textbf{44)} Conferi o enunciado na p. 15 do PDF (questão 44, impresso como p. 13) e a transcrição está fiel. Refiz a questão lacuna por lacuna. (i) A rede nervosa difusa surge nos cnidários. (ii) Os cordões nervosos longitudinais e os gânglios cerebroides anteriores (cefalização, sistema em escada) surgem nos platelmintos. (iii) O potencial de ação ocorre na membrana plasmática. (iv) Na despolarização, o Na+ entra. (v) Na repolarização, o K+ sai. Só a alternativa B tem as cinco lacunas certas. As outras falham porque põem equinodermos em (i) ou (ii), citam o envoltório nuclear ou invertem os íons. A resolução e o gabarito coincidem e estão certos.%
\end{itemize}%
\paragraph{Pontos de dúvida}%
\begin{itemize}%
\item%
\textit{(statement)} O enunciado diz que a saída de K+ vai 'retornando o gradiente eletroquímico ao estado de repouso'. A rigor, o efluxo de K+ restabelece o potencial de membrana. As concentrações iônicas são restabelecidas pela bomba Na+/K+{-}ATPase. A resolução já explica essa distinção, e ela não muda a alternativa.%
\item%
\textit{(solution)} A resolução diz que os platelmintos são 'os primeiros animais com simetria bilateral'. Essa é a simplificação usual dos livros didáticos. Em filogenias modernas, os acelomorfos são o grupo{-}irmão dos demais bilatérios. Isso não afeta a resposta.%
\end{itemize}%
\section*{Questão 45}%
\begin{longtable}{@{}p{0.06\textwidth} p{0.33\textwidth} p{0.33\textwidth} p{0.1\textwidth} p{0.08\textwidth}@{}}\toprule Item & Resolução & Gabarito & Veredito & Certo \\ \midrule
45 & (C) Há fototropismo positivo na parte aérea, e a concentração de auxina é maior na seta 2 (lado sombreado) do que na seta 3 (lado iluminado). & C & igual & ambos \\
\bottomrule\end{longtable}%
\begin{itemize}%
\item%
\textbf{45)} Conferi a figura no PDF (p. 16 do arquivo). O Sol está no alto, à esquerda, e o caule se curva para ele. As setas 1 e 3 ficam no lado esquerdo do caule, o lado iluminado. A seta 2 fica no lado direito do alto do caule, o lado sombreado. As setas 4 e 5 marcam os lados da raiz, e a 6 marca a ponta da raiz. Pela hipótese de Cholodny--Went, a auxina migra para o lado sombreado, que se alonga mais, e o caule se curva para a luz. Então {[}auxina{]} é maior em 2 do que em 3, e a (C) está correta. A seta 2 também fica mais perto do ápice, onde a auxina é produzida, o que só reforça a desigualdade. Por que as outras estão erradas: (a) põe o aumento de auxina no lado iluminado (seta 1). (b) e (d) atribuem o gravitropismo à citocinina, mas o hormônio responsável é a auxina. A resolução e o gabarito batem e os dois estão certos.%
\end{itemize}%
\paragraph{Pontos de dúvida}%
\begin{itemize}%
\item%
\textit{(solution)} Ao descartar a (d), a resolução diz que 'a raiz da figura cresce verticalmente'. Na figura, a raiz principal tem uma leve curvatura. Isso não muda o resultado: a (d) é falsa porque fala em citocinina, não por causa da forma da raiz.%
\item%
\textit{(solution)} Na (a), a resolução trata a seta 1 como o lado iluminado da região apical. Na figura, a seta 1 aponta para o lado esquerdo do caule, logo abaixo do ápice, e não para o meristema em si. A leitura é razoável e não altera a conclusão.%
\end{itemize}%
\section*{Questão 46}%
\begin{longtable}{@{}p{0.06\textwidth} p{0.33\textwidth} p{0.33\textwidth} p{0.1\textwidth} p{0.08\textwidth}@{}}\toprule Item & Resolução & Gabarito & Veredito & Certo \\ \midrule
46 & (B) pirâmides de biomassa; produtores e consumidores primários impactados por desequilíbrio, teia ainda se reestruturando em t2 & B & igual & ambos \\
\bottomrule\end{longtable}%
\begin{itemize}%
\item%
\textbf{46)} Conferi a figura no original (p. 16 do PDF, p. 14 da prova). Em t1 a pirâmide é regular. Em t2 os produtores (branco) e os consumidores primários (cinza{-}claro) encolhem muito, enquanto os consumidores secundários (cinza{-}escuro) e terciários (preto) ficam com largura praticamente igual. Com isso, t2 fica parcialmente invertida. Pirâmide de energia não se inverte, o que elimina a e c. Sobra biomassa. Em t2 os níveis superiores ainda não acompanharam a queda dos níveis que os sustentam, então é um estado transitório e não um equilíbrio restabelecido. Isso elimina a d, que também menciona só os produtores. A resolução está correta e bate com o gabarito.%
\end{itemize}%
\paragraph{Pontos de dúvida}%
\begin{itemize}%
\item%
\textit{(solution)} A resolução atribui a diminuição do fluxo entre níveis à 'conservação da energia'. O mais preciso seria dizer que a perda de energia como calor (respiração, 2ª lei da termodinâmica) e a parte não assimilada fazem o fluxo cair de um nível para o outro. A conclusão continua certa: a pirâmide de energia nunca se inverte.%
\item%
\textit{(statement)} A leitura das larguras depende da figura. Pela comparação visual, os níveis cinza{-}escuro e preto têm praticamente a mesma largura em t1 e t2. Isso apoia a ideia de que os consumidores do topo ainda não reagiram à queda dos níveis de baixo, mas é uma inferência visual, não um dado quantitativo.%
\end{itemize}%
\section*{Questão 47}%
\begin{longtable}{@{}p{0.06\textwidth} p{0.33\textwidth} p{0.33\textwidth} p{0.1\textwidth} p{0.08\textwidth}@{}}\toprule Item & Resolução & Gabarito & Veredito & Certo \\ \midrule
Q47 & D & D & igual & ambos \\
\bottomrule\end{longtable}%
\begin{itemize}%
\item%
\textbf{Q47)} Conferi o enunciado na página 16 do PDF (página 14 do caderno) e a transcrição está fiel. A alternativa D é a correta. A AIDS é o estágio mais avançado da infecção pelo HIV, e a terapia antirretroviral, ao deixar a carga viral indetectável, impede a progressão para a doença. As outras três estão erradas. Em A, o sistema imune não elimina o HIV, e sem tratamento a maioria dos infectados evolui para AIDS. Em B, os antirretrovirais previnem, sim, a transmissão (indetectável = intransmissível) e a progressão. Em C, a relação causal está invertida: a imunodeficiência deixa a pessoa suscetível às doenças oportunistas, e não existem 'doenças que causam a AIDS'. A análise das alternativas na resolução é correta e coerente.%
\end{itemize}%
\section*{Questão 48}%
\begin{longtable}{@{}p{0.06\textwidth} p{0.33\textwidth} p{0.33\textwidth} p{0.1\textwidth} p{0.08\textwidth}@{}}\toprule Item & Resolução & Gabarito & Veredito & Certo \\ \midrule
48 & (C) substância branca, na região periférica, formada por fibras nervosas que conduzem impulsos ascendentes ao encéfalo e descendentes gerados por ele & C & igual & ambos \\
\bottomrule\end{longtable}%
\begin{itemize}%
\item%
\textbf{48)} Conferi o enunciado na página 16 do PDF (página 14 do caderno) e ele bate com a transcrição. No corte transversal da medula, a substância cinzenta fica no centro, em forma de H, e reúne os corpos celulares, os dendritos, os interneurônios e as sinapses. A substância branca fica na periferia e é formada por axônios mielinizados, organizados em tratos ascendentes (sensitivos) e descendentes (motores). Só a alternativa (c) junta a localização certa (periférica) com a composição certa (fibras nervosas). A (a) põe a periferia na substância cinzenta. A (b) diz que a substância branca é central e tem corpos celulares. A (d) descreve corretamente a substância cinzenta, mas ali se concentram corpos celulares, e não axônios. A resolução analisa as quatro alternativas sem erro conceitual, e a ligação que faz com a perda motora e sensorial abaixo da lesão é coerente.%
\end{itemize}%
\paragraph{Pontos de dúvida}%
\begin{itemize}%
\item%
\textit{(statement)} O enunciado usa 'concentram{-}se', o que trata como simplificação aceita o fato de que a substância cinzenta também tem axônios (curtos e pouco mielinizados, de interneurônios, além das porções iniciais dos axônios dos motoneurônios). Isso não cria ambiguidade entre as alternativas: a (c) é a única coerente com a anatomia.%
\end{itemize}%
\section*{Questão 49}%
\begin{longtable}{@{}p{0.06\textwidth} p{0.33\textwidth} p{0.33\textwidth} p{0.1\textwidth} p{0.08\textwidth}@{}}\toprule Item & Resolução & Gabarito & Veredito & Certo \\ \midrule
Q49 & D & D & igual & ambos \\
\bottomrule\end{longtable}%
\begin{itemize}%
\item%
\textbf{Q49)} Conferi a transcrição com o original (páginas 14--15 impressas no caderno, que são as páginas 16--17 do PDF) e ela está fiel. O texto diz que a terceira dimensão é uma "tendência", que já se percebe na humanidade misturada e concentrada nas grandes cidades. Também a apresenta, ao lado da fábula e da perversidade, como uma das três dimensões que é preciso considerar juntas. Isso confirma a (d): unidade, totalidade aberta e transformações já em curso nas grandes cidades. As outras caem pelos motivos que a resolução aponta. A (a) repete a fábula da "aldeia global". A (b) diz que as dimensões estão dissociadas e trata a terceira como fábula. A (c) fica só na perversidade e fala em "homogeneização", quando Santos descreve uma modernização seletiva que aprofunda as desigualdades.%
\end{itemize}%
\section*{Questão 50}%
\begin{longtable}{@{}p{0.06\textwidth} p{0.33\textwidth} p{0.33\textwidth} p{0.1\textwidth} p{0.08\textwidth}@{}}\toprule Item & Resolução & Gabarito & Veredito & Certo \\ \midrule
50 & (B) EUA e China, economias interdependentes e concorrentes que dominam a infraestrutura de IA e exercem hegemonia num mundo multipolar em transformação acelerada, com dilemas ético e existencial candentes. & B & igual & ambos \\
\bottomrule\end{longtable}%
\begin{itemize}%
\item%
\textbf{50)} A transcrição confere com a página 17 do PDF (página 15 da prova): o texto e as quatro alternativas estão iguais. A resolução marca B, a mesma alternativa do gabarito oficial. A justificativa se apoia no texto. A China alerta os EUA, e o texto fala no 'intenso debate' ético dentro dos EUA. Isso derruba a ideia de ordem unipolar em (a) e a de ausência de critérios éticos em (d). Também não há base para dizer que a China seguiu tratados de não proliferação, como afirma (c). A rejeição das alternativas é coerente: (a) fala em unipolaridade, (c) é uma leitura idealizada da China e (d) fala em bipolaridade, economias 'independentes' e ausência genérica de ética. O gabarito está correto.%
\end{itemize}%
\paragraph{Pontos de dúvida}%
\begin{itemize}%
\item%
\textit{(statement)} A alternativa B, que é a oficial, junta 'exercem hegemonia' com 'mundo multipolar'. As duas ideias ficam em certa tensão conceitual, porque hegemonia compartilhada por dois países lembra bipolaridade. Ainda assim, B é a única alternativa sem erro claro, então isso não muda a resposta. A resolução trata esse ponto bem: diz que a hegemonia tecnológica sino{-}americana se exerce dentro de um quadro multipolar.%
\item%
\textit{(solution)} A resolução traz afirmações de contexto que o texto não contém. São elas: 'terras raras' como parte da infraestrutura de IA e a ideia de que 'não existe tratado internacional consolidado sobre armas autônomas com IA'. As duas são razoáveis e não afetam a conclusão, mas são conhecimento externo, não algo tirado do enunciado.%
\end{itemize}%
\section*{Questão 51}%
\begin{longtable}{@{}p{0.06\textwidth} p{0.33\textwidth} p{0.33\textwidth} p{0.1\textwidth} p{0.08\textwidth}@{}}\toprule Item & Resolução & Gabarito & Veredito & Certo \\ \midrule
51 & (A) adaptação das usinas para processar cana e milho (usinas flex), com maior flexibilidade na produção do biocombustível & A & igual & ambos \\
\bottomrule\end{longtable}%
\begin{itemize}%
\item%
\textbf{51)} Conferi o enunciado na página 17 do PDF (página 15 do caderno). A transcrição está fiel: texto, gráfico com cana caindo de \textasciitilde{}93\% para \textasciitilde{}52\% e milho subindo de \textasciitilde{}7\% para \textasciitilde{}48\% entre 2017 e 2024, e as quatro alternativas. A alternativa (a) descreve as usinas flex, que moem cana na safra e usam milho armazenado na entressafra, o que dá mais flexibilidade. As justificativas para descartar as outras estão corretas. (b) erra porque a cana é perecível e precisa ficar perto da usina, enquanto o milho é grão seco, pode ser estocado em silos e transportado de longe, então não há a 'correspondência na estocagem' que a alternativa afirma. (c) erra porque a cana é semiperene: o canavial continua no campo durante a entressafra, que é só o período sem colheita. O milho para etanol em Mato Grosso vem sobretudo da safrinha plantada depois da soja. (d) erra porque a colheita de milho nas grandes propriedades é amplamente mecanizada. A resolução e o gabarito oficial marcam (A), e os dois estão certos.%
\end{itemize}%
\paragraph{Pontos de dúvida}%
\begin{itemize}%
\item%
\textit{(solution)} Ao descartar a (c), a resolução diz que 'não existe rotação com milho' em área de cana. Isso é um pouco absoluto. Na reforma do canavial, a cada 5 a 6 anos, às vezes se faz rotação, quase sempre com leguminosas como soja, amendoim ou crotalária. Mas a alternativa fala em plantar milho na entressafra da cana, e nesse período a cana continua no campo. Por isso a (c) continua errada e a nuance não muda a resposta.%
\item%
\textit{(solution)} O prazo de '24 a 48 horas' entre o corte e a moagem da cana é uma estimativa usual. As fontes variam entre cerca de 24 e 72 horas. É só um detalhe ilustrativo e não interfere na análise.%
\item%
\textit{(solution)} Os percentuais citados (\textasciitilde{}93\%/7\% em 2017 e \textasciitilde{}52\%/48\% em 2024) foram lidos no gráfico de barras e são aproximados. Batem com a figura do original, e a resposta não depende dos valores exatos.%
\end{itemize}%
\section*{Questão 52}%
\begin{longtable}{@{}p{0.06\textwidth} p{0.33\textwidth} p{0.33\textwidth} p{0.1\textwidth} p{0.08\textwidth}@{}}\toprule Item & Resolução & Gabarito & Veredito & Certo \\ \midrule
52 & A & A & igual & ambos \\
\bottomrule\end{longtable}%
\begin{itemize}%
\item%
\textbf{52)} Conferi o enunciado na página 18 do PDF (página 16 da prova), e a transcrição está fiel. A resolução e o gabarito oficial da Comvest marcam a alternativa (A). Recalculei as dez taxas da tabela com (P2022-P2010)/P2010$\cdot$100\%, e todas batem com os valores impressos (84,31; 82,27; \ldots{}; 50,01). A explicação da resolução está correta. Ela separa dois grupos. No primeiro estão as cidades de fronteira econômica: agrícola (Luís Eduardo Magalhães, Sinop, Sorriso), mineral (Parauapebas, em Carajás), do petróleo (Maricá, com royalties do pré{-}sal da Bacia de Santos) e do turismo (Camboriú). No segundo estão as cidades que crescem pela dispersão metropolitana com expansão imobiliária: Senador Canedo (RM de Goiânia), Fazenda Rio Grande (RM de Curitiba), Palhoça (Grande Florianópolis) e São José de Ribamar (Grande São Luís). A alternativa (A) é a única que junta esses dois vetores. A resolução descarta bem as outras: (b) porque o crescimento não vem da indústria de base nem de polos portuários corporativos; (c) porque a premissa é falsa, já que a violência se interiorizou; (d) porque os programas de moradia não transferem população das metrópoles para o interior.%
\end{itemize}%
\paragraph{Pontos de dúvida}%
\begin{itemize}%
\item%
\textit{(solution)} A resolução põe cada município em uma única categoria, usando conhecimento que não está no enunciado. Alguns casos se encaixam em mais de um vetor. Maricá também cresce pela expansão a partir da Região Metropolitana do Rio. São José de Ribamar e Palhoça também têm atividade litorânea e turística. Isso não muda a escolha da (A), porque a alternativa cita os dois tipos de vetor juntos.%
\end{itemize}%
\section*{Questão 53}%
\begin{longtable}{@{}p{0.06\textwidth} p{0.33\textwidth} p{0.33\textwidth} p{0.1\textwidth} p{0.08\textwidth}@{}}\toprule Item & Resolução & Gabarito & Veredito & Certo \\ \midrule
Q53 & (a) Os desertos ficam majoritariamente na zona relativamente natural (A), e as savanas se distribuem de forma mais equilibrada entre as zonas A, B e C. & A & igual & ambos \\
\bottomrule\end{longtable}%
\begin{itemize}%
\item%
\textbf{Q53)} Conferi o mapa no PDF (página 18, impressa como 16). O deserto, em amarelo{-}claro, tem um único rótulo A, que cobre o grande bloco central e ocidental. Os rótulos C aparecem só nas bordas: litoral oeste, Grande Baía Australiana e o lobo leste. A savana, em laranja, tem rótulos A (Kimberley, Top End, mancha central, interior de Queensland, faixa do sul da Austrália Ocidental) e C (interior do Top End, interior do norte de Queensland e, por linhas de chamada, o litoral norte). Também tem B na parte de Queensland. Isso confirma (a). A (b) inverte os biomas. A (c) é falsa, porque o interior desértico é quase todo zona A e a agricultura e as cidades estão na periferia úmida. A (d) é falsa, porque há rótulos C sobre a savana, e a justificativa ainda é determinista. A resolução e o gabarito coincidem e estão certos.%
\end{itemize}%
\paragraph{Pontos de dúvida}%
\begin{itemize}%
\item%
\textit{(solution)} Alguns detalhes da leitura do mapa dependem de rótulos pequenos e de linhas de chamada. São eles: a contagem de exatamente quatro rótulos C sobre o deserto, a afirmação de que não há nenhum B sobre ele e a atribuição dos B de Queensland ao interior da savana ou às faixas de floresta. Por exemplo, o B no extremo norte, perto do Cabo York, e os B/C no sul da Austrália, perto dos Montes Flinders, ficam em fronteiras entre biomas. Uma leitura diferente desses trechos não muda a conclusão, porque o deserto continua majoritariamente A e a savana continua com as três zonas.%
\item%
\textit{(statement)} A alternativa (a) usa termos comparativos e vagos, como 'majoritariamente' e 'distribuição mais equilibrada'. A zona A ainda é a maior dentro da savana. A alternativa só se sustenta como comparação com o deserto, e não como uma divisão em partes iguais. As outras alternativas são claramente falsas, então não há ambiguidade real.%
\end{itemize}%
\section*{Questão 54}%
\begin{longtable}{@{}p{0.06\textwidth} p{0.33\textwidth} p{0.33\textwidth} p{0.1\textwidth} p{0.08\textwidth}@{}}\toprule Item & Resolução & Gabarito & Veredito & Certo \\ \midrule
54 & (C) & C & igual & ambos \\
\bottomrule\end{longtable}%
\begin{itemize}%
\item%
\textbf{54)} Conferi os dois mapas no original (pág. 19 do PDF, renderizada em 300 dpi). A leitura da resolução está certa. Mapa 1 (P): BH fica na mancha de 1400--1600, mais escura que o entorno; Rio, São Paulo e Vitória ficam na classe 1600--1800. Mapa 2 (ETP): BH em 1400--1600; Rio com o oeste em 1400--1600 e o leste em 1600--1800; São Paulo com o oeste em 1400--1600 e o leste em <1400; Vitória em 1600--1800. Assim, só em SP a classe de P fica toda acima da de ETP, e a premissa de (c) se confirma. A conclusão de (c), que pede monitoramento por causa da sazonalidade, do armazenamento e da demanda, é coerente com o enunciado. (a) e (b) caem pelas conclusões ('disponibilidade ao longo do ano', 'assegura', 'dispensando alertas'). (d) cai porque P e ETP de Vitória estão na mesma classe. As faixas de excedente também conferem: (0, 400) mm no oeste e >200 mm no leste.%
\end{itemize}%
\paragraph{Pontos de dúvida}%
\begin{itemize}%
\item%
\textit{(statement)} A transcrição do enunciado põe a RM de Vitória em área de 1200--1400 mm de precipitação. No mapa original ampliado, a RM está dentro da mancha de 1600--1800 mm, que é a leitura usada na resolução. Para a alternativa (d), as duas leituras dão o mesmo resultado: com P em 1200--1400 e ETP em 1600--1800, (d) teria premissa verdadeira, mas a resposta continuaria (c), porque a (c) é a única com premissa e conclusão corretas sem ambiguidade. Mesmo assim, vale confirmar a leitura no original.%
\item%
\textit{(solution)} A resolução traz contexto externo (Cantareira, crise de 2014--2015, \textasciitilde{}20 milhões de habitantes) que não está no enunciado. Não muda a resposta, é só ilustração.%
\item%
\textit{(statement)} A RM do Rio cai em duas classes de ETP: 1400--1600 a oeste e 1600--1800 a leste. Por isso, a afirmação de 'faixas distintas' na (b) é só parcialmente verdadeira. A (b) cai de qualquer forma pela conclusão absoluta.%
\end{itemize}%
\section*{Questão 55}%
\begin{longtable}{@{}p{0.06\textwidth} p{0.33\textwidth} p{0.33\textwidth} p{0.1\textwidth} p{0.08\textwidth}@{}}\toprule Item & Resolução & Gabarito & Veredito & Certo \\ \midrule
Questão 55 & (A): o logotipo identifica e simboliza a universidade. Ele lembra a organização espacial do campus, mas é uma representação gráfica simplificada ligada à identidade institucional, e não um mapa. & A & igual & ambos \\
\bottomrule\end{longtable}%
\begin{itemize}%
\item%
\textbf{Questão 55)} Comparei a transcrição com o original (p. 19 do PDF, numerada 17 na prova), e o enunciado e as alternativas são os mesmos. A resolução separa bem a representação simbólica, que é o logotipo, da representação cartográfica. O logotipo não tem escala, orientação nem legenda, e o desenho não corresponde ponto a ponto às vias e aos edifícios. Por isso as alternativas b e d estão erradas ao dizer que ele é um mapa, e a c está errada ao dizer que ele reproduz a distribuição dos edifícios e das vias e funciona como representação cartográfica. A alternativa a é a única que reconhece a função identitária do logotipo, admite o diálogo visual com o campus e o chama de representação gráfica simplificada. A resolução e o gabarito oficial concordam: a resposta é A.%
\end{itemize}%
\paragraph{Pontos de dúvida}%
\begin{itemize}%
\item%
\textit{(solution)} A resolução diz que os segmentos pretos não correspondem um a um aos edifícios e vias reais e que os três círculos vermelhos não são lugares localizáveis. Isso vem da comparação visual com a foto aérea e de conhecimento geral sobre a marca, e o enunciado não diz isso. Mesmo assim, a conclusão não depende desse detalhe, porque a falta de escala, orientação e legenda já basta para mostrar que o logotipo não é um mapa.%
\end{itemize}%
\section*{Questão 56}%
\begin{longtable}{@{}p{0.06\textwidth} p{0.33\textwidth} p{0.33\textwidth} p{0.1\textwidth} p{0.08\textwidth}@{}}\toprule Item & Resolução & Gabarito & Veredito & Certo \\ \midrule
única (alternativa) & B & B & igual & ambos \\
\bottomrule\end{longtable}%
\begin{itemize}%
\item%
\textbf{única (alternativa))} Reli a figura no PDF (página 20) e ela bate com a transcrição: All 32, Some 51, None 16, e no subgrupo 'some' 97/52/44/15/14/9/5. A alternativa B se apoia nos dados: entre os que apoiam deportação condicionada, só 14\% deportariam pais de crianças nascidas nos EUA e só 5\% deportariam casados com cidadão americano, contra 97\% para crimes violentos. Refiz as contas que descartam as demais. (a) 32\% fica abaixo de 50\%. (c) 0,52$\cdot$51\% $\approx$ 26,5\%; se somar os 32\% do grupo 'all', dá cerca de 58,5\%; nenhum dos dois é cerca de um terço. (d) 0,97$\cdot$51\% $\approx$ 49,5\%; com o grupo 'all', dá cerca de 81,5\%; além disso, 16\% rejeitam qualquer deportação, então não é a 'quase totalidade'. As contas da resolução estão certas e a conclusão coincide com o gabarito.%
\end{itemize}%
\paragraph{Pontos de dúvida}%
\begin{itemize}%
\item%
\textit{(solution)} Na alternativa (d), a resolução diz que apoiam a deportação de autores de crimes violentos 'no máximo' 32\% + 0,97$\cdot$51\% $\approx$ 81,5\%. Seria mais exato dizer 'aproximadamente'. Quem quer deportar todos (all) também apoia, logicamente, deportar quem cometeu crime violento, então esse valor é uma estimativa e não um teto. Isso não muda a conclusão: 16\% rejeitam qualquer deportação.%
\item%
\textit{(statement)} Em (b), a comparação 'mais tolerantes' não diz com o quê se compara: com os critérios de crime dentro do mesmo subgrupo ou com o grupo que defende deportar todos. As duas leituras levam a B, então a ambiguidade não afeta a resposta.%
\end{itemize}%
\section*{Questão 57}%
\begin{longtable}{@{}p{0.06\textwidth} p{0.33\textwidth} p{0.33\textwidth} p{0.1\textwidth} p{0.08\textwidth}@{}}\toprule Item & Resolução & Gabarito & Veredito & Certo \\ \midrule
Questão 57 & (A) natureza transgressora da arte; apenas Morrison cita quais tipos de pessoas têm atacado artistas & A & igual & ambos \\
\bottomrule\end{longtable}%
\begin{itemize}%
\item%
\textbf{Questão 57)} Conferi a transcrição com o original (p. 20 do PDF, p. 18 da prova) e ela confere. A resposta A está certa. A arte é chamada de 'revolutionary' por Moura e de 'dangerous' por Morrison, então a ideia comum é a natureza transgressora da arte. Só Morrison diz quem ataca os artistas: 'dictators and people in office and people who want to control and deceive'. Moura fala apenas em 'these people'/'they', sem dizer quem são. A B inverte o que Moura diz: 'They're not afraid of what I'm saying {[}...{]} They're afraid of what we do'. A C depende do falso cognato deceive = enganar, e não 'decepcionar'. A D está errada na primeira parte, porque nenhum dos dois fala em uma luta contra a punição de artistas, ainda que seja verdade que só Moura cita o cinema ('our films'). A análise da resolução está correta e completa.%
\end{itemize}%
\section*{Questão 58}%
\begin{longtable}{@{}p{0.06\textwidth} p{0.33\textwidth} p{0.33\textwidth} p{0.1\textwidth} p{0.08\textwidth}@{}}\toprule Item & Resolução & Gabarito & Veredito & Certo \\ \midrule
58 & (C): frutos vermelhos em forma de gota, cobertos de cerdas flexíveis (urucum, Bixa orellana) & C & igual & ambos \\
\bottomrule\end{longtable}%
\begin{itemize}%
\item%
\textbf{58)} Conferi na página 20 do PDF original. A ilustração c) mostra dois frutos vermelhos em forma de gota, cobertos de cerdas flexíveis, com aspecto peludo. É a cápsula do urucum, que bate com 'teardrop{-}shaped capsules', 'malleable bristles' e 'fuzzy appearance'. O nome vem do tupi e alude ao corante vermelho, e as sementes servem para pintura corporal e repelente. As outras ilustrações não servem: a) é um fruto liso e sulcado, tipo cacau; b) são frutos redondos e alaranjados, tipo tucumã/pupunha; d) são frutos abertos com semente preta em arilo branco, tipo guaraná. A análise das alternativas na resolução está correta e coincide com o gabarito oficial.%
\end{itemize}%
\section*{Questão 59}%
\begin{longtable}{@{}p{0.06\textwidth} p{0.33\textwidth} p{0.33\textwidth} p{0.1\textwidth} p{0.08\textwidth}@{}}\toprule Item & Resolução & Gabarito & Veredito & Certo \\ \midrule
59 & (B) esclarece a função de algumas gírias na demarcação de território e identidade cultural & B & igual & ambos \\
\bottomrule\end{longtable}%
\begin{itemize}%
\item%
\textbf{59)} Conferi a transcrição com o original (págs. 18--19 do caderno) e ela está fiel. O texto diz que as gírias "signaled in{-}group status" e que grupos marginalizados criam gírias "as a tool of power to build community and shared identity". Isso corresponde à alternativa (b). As outras alternativas estão erradas. (a) O autor não defende as redes sociais: ele descreve a difusão por apropriação e a perda da exclusividade. (c) O texto não fala em direito de recusar a apropriação; o autor prevê que ela continua ("it's going to be taken too", "never{-}ending cycle"). (d) Inverte os papéis: a exclusividade é dos grupos marginalizados, e os grupos dominantes apenas se apropriam das gírias pelo "perceived coolness".%
\end{itemize}%
\paragraph{Pontos de dúvida}%
\begin{itemize}%
\item%
\textit{(solution)} Na análise da alternativa (d) há um erro de digitação: em \textbackslash{}emph\{marginalizados*\} sobrou um asterisco. Não muda o sentido.%
\item%
\textit{(solution)} A resolução lê "território" como pertencimento ao grupo (in{-}group status). O texto não fala em território físico; trata{-}se de território simbólico. A leitura é razoável, coincide com o gabarito e não afeta a resposta.%
\end{itemize}%
\section*{Questão 60}%
\begin{longtable}{@{}p{0.06\textwidth} p{0.33\textwidth} p{0.33\textwidth} p{0.1\textwidth} p{0.08\textwidth}@{}}\toprule Item & Resolução & Gabarito & Veredito & Certo \\ \midrule
Questão 60 & C & C & igual & ambos \\
\bottomrule\end{longtable}%
\begin{itemize}%
\item%
\textbf{Questão 60)} Conferi a transcrição contra a página 21 do PDF (impressa como p. 19) e ela está correta. A resposta é (c) porque a primeira frase diz que os AI companion products "can commodify intimacy through emotionally manipulative design". Na alternativa, "a forma como são desenvolvidos" corresponde a design, "tem o potencial" a can e "transformar a intimidade em mercadoria" a commodify intimacy. As outras alternativas estão erradas. (a) Economia política e estudos culturais são a lente de análise do artigo, não o tema dos estereótipos. (b) O texto diz o contrário de "apesar da reciprocidade", pois a estrutura do produto "foreclose{[}s{]} genuinely reciprocal relationships". (d) Os produtos recorrem a estereótipos de servidão e docilidade ("draw on", "repackage") em vez de combatê{-}los. A resolução analisa as quatro alternativas corretamente.%
\end{itemize}%
\paragraph{Pontos de dúvida}%
\begin{itemize}%
\item%
\textit{(statement)} As alternativas falam em "assistentes de IA", mas o texto trata de "AI companions" (companheiros de IA). A diferença é só de tradução e não afeta a escolha da alternativa (c).%
\end{itemize}%
\section*{Questão 61}%
\begin{longtable}{@{}p{0.06\textwidth} p{0.33\textwidth} p{0.33\textwidth} p{0.1\textwidth} p{0.08\textwidth}@{}}\toprule Item & Resolução & Gabarito & Veredito & Certo \\ \midrule
Q61 & (D) & D & igual & ambos \\
\bottomrule\end{longtable}%
\begin{itemize}%
\item%
\textbf{Q61)} Conferi o enunciado na página 21 do PDF e ele bate com a transcrição. A resposta D está correta. O poema contrapõe o passado de violência ao presente de resistência. O trauma aparece em formas de passado e particípio: "who's carried the cross", "Gold stripped", "names erased", "Bones buried", "those who burned / bled / built and were buried". A resistência aparece no presente simples: "Yet she breathes / she births / she dances", "I carry her", "I am her memory", "I wear it for her", "It pulses", "And still we rise". As outras alternativas são falsas. Em (a), a coroa simboliza dignidade ("Not gold but grace... A crown that cannot be taken"), e quem simboliza o sofrimento é a cruz. Em (b), o presente simples marca resistência, não a persistência da opressão. Em (c), há personificação, mas ela mostra um corpo vivo e materno, que não perdeu as raízes. A análise da resolução está correta e bem fundamentada.%
\end{itemize}%
\paragraph{Pontos de dúvida}%
\begin{itemize}%
\item%
\textit{(statement)} Na alternativa (a), o enunciado oficial escreve "sofrimento infringido". O correto seria "infligido" (de infligir). É erro de redação da prova e não afeta a resposta.%
\item%
\textit{(solution)} A alternância descrita em (d) não vale para todos os verbos. Alguns particípios estão do lado positivo: "A continent crowned by salt and sun" e "we rise crowned". Além disso, "built" não indica trauma por si só. A resolução lista "built" entre os verbos de passado sem fazer essa ressalva. Mesmo assim, (d) continua sendo claramente a melhor alternativa.%
\end{itemize}%
\section*{Questão 62}%
\begin{longtable}{@{}p{0.06\textwidth} p{0.33\textwidth} p{0.33\textwidth} p{0.1\textwidth} p{0.08\textwidth}@{}}\toprule Item & Resolução & Gabarito & Veredito & Certo \\ \midrule
62 & D & D & igual & ambos \\
\bottomrule\end{longtable}%
\begin{itemize}%
\item%
\textbf{62)} Conferi a transcrição com o original (págs. 19--20 da prova, 21--22 do PDF) e ela está fiel. A alternativa (d) é a correta. O texto diz que "all sentient beings, human or nonhuman" têm pelo menos dois interesses, o de não sofrer e o de não morrer, e que "the use of animals as property ... implicates at least two interests that animals have". Logo, esses interesses são comuns a humanos e animais, e o uso de animais como propriedade os fere. As outras alternativas estão erradas. A (a) contradiz o texto, que afirma "There is absolutely no doubt that the animals we routinely exploit are sentient". A (b) fala em criminalizar, o que não aparece no texto: a obrigação é deixar claro que a exploração é moralmente injustificável "however 'humane'". A (c) inverte o texto, segundo o qual julgar pessoas "is never a good idea" e o essencial é julgar ações. A análise da resolução está correta em todas as alternativas.%
\end{itemize}%
\paragraph{Pontos de dúvida}%
\begin{itemize}%
\item%
\textit{(solution)} Ao refutar (a), a resolução cita a alternativa como "mesmo que alguns deles não sejam capazes de sentir", mas o texto original diz "sentir emoções". Para os autores, senciência é a capacidade de sofrer, o que não é exatamente "sentir emoções". Isso é só imprecisão de redação: a refutação continua válida, porque o texto afirma sem ressalvas que todos os animais explorados são sencientes.%
\end{itemize}%
\section*{Questão 63}%
\begin{longtable}{@{}p{0.06\textwidth} p{0.33\textwidth} p{0.33\textwidth} p{0.1\textwidth} p{0.08\textwidth}@{}}\toprule Item & Resolução & Gabarito & Veredito & Certo \\ \midrule
Q63 & (a) d $\approx$ 1,2$\times$10\^{}25 m & A & igual & ambos \\
\bottomrule\end{longtable}%
\begin{itemize}%
\item%
\textbf{Q63)} Conferi o enunciado no PDF original (pág. 22): c = 3,0$\times$10\^{}8 m/s, 1 ano $\approx$ 3,2$\times$10\^{}7 s, tempo de 1,3 bilhão de anos. Refazendo a conta: $\Delta$t = 1,3$\times$10\^{}9 $\times$ 3,2$\times$10\^{}7 = 4,16$\times$10\^{}16 s. Daí d = 3,0$\times$10\^{}8 $\times$ 4,16$\times$10\^{}16 = 1,248$\times$10\^{}25 m $\approx$ 1,2$\times$10\^{}25 m, que é a alternativa (a). A origem dos distratores também confere: (b) 3,0$\times$10\^{}8$\cdot$1,3$\times$10\^{}9 = 3,9$\times$10\^{}17, sem converter anos em segundos; (c) bilhão tomado como 10\^{}6, com a conversão feita, dá 1,248$\times$10\^{}22; (d) bilhão como 10\^{}6 e sem conversão dá 3,9$\times$10\^{}14. A resolução e o gabarito estão corretos.%
\end{itemize}%
\paragraph{Pontos de dúvida}%
\begin{itemize}%
\item%
\textit{(solution)} A resolução trata a distância como d = c$\cdot$$\Delta$t e despreza a expansão do Universo. É o modelo que a questão pede e a hipótese foi declarada, mas a distância cosmológica real seria outra. Isso não muda a alternativa.%
\item%
\textit{(solution)} Há uma diferença de arredondamento: o corpo da resolução dá $\approx$1,2$\times$10\^{}25 m e a resposta final declarada dá $\approx$1,25$\times$10\^{}25 m (o valor exato é 1,248$\times$10\^{}25 m). Os dois levam à alternativa (a).%
\end{itemize}%
\section*{Questão 64}%
\begin{longtable}{@{}p{0.06\textwidth} p{0.33\textwidth} p{0.33\textwidth} p{0.1\textwidth} p{0.08\textwidth}@{}}\toprule Item & Resolução & Gabarito & Veredito & Certo \\ \midrule
64 & (B): independe de m, é diretamente proporcional a M e inversamente proporcional ao quadrado do raio da Terra (g = GM/R$^2$) & B & igual & ambos \\
\bottomrule\end{longtable}%
\begin{itemize}%
\item%
\textbf{64)} Conferi o enunciado na página 22 do PDF (página 20 do caderno) e a transcrição está fiel. A massa m se cancela: pela gravitação, F = GMm/R$^2$, e pela 2ª lei de Newton, g = F/m = GM/R$^2$. Assim, g é independente de m, diretamente proporcional a M e inversamente proporcional a R$^2$, que é a alternativa B. A alternativa c descreve a dependência da força peso, não da aceleração; as alternativas a e d trazem g diretamente proporcional a R$^2$, o que é falso. A conta de ordem de grandeza da resolução está certa: 6,67e{-}11 $\times$ 5,97e24 / (6,37e6)$^2$ $\approx$ 9,8 m/s$^2$. A resolução e o gabarito coincidem e estão corretos.%
\end{itemize}%
\section*{Questão 65}%
\begin{longtable}{@{}p{0.06\textwidth} p{0.33\textwidth} p{0.33\textwidth} p{0.1\textwidth} p{0.08\textwidth}@{}}\toprule Item & Resolução & Gabarito & Veredito & Certo \\ \midrule
Q65 & (D) 240 N: F = T1 = 2T = 2mg & D & igual & ambos \\
\bottomrule\end{longtable}%
\begin{itemize}%
\item%
\textbf{Q65)} Ampliei a figura original a 600 dpi para conferir a montagem, que confirma a leitura da resolução. O fio da manopla passa pelas duas polias fixas e desce até o EIXO da polia móvel. O segundo fio passa por cima da polia móvel: o ramo da esquerda está preso à base e o da direita sustenta o bloco. O rótulo "Fios" é só uma seta de legenda, não é um fio. Equilíbrio do bloco: T = mg = 120 N. Na polia móvel, sem massa, os dois ramos puxam para baixo com T cada, então a tração no eixo é T1 = 2T = 240 N. As polias fixas ideais não alteram a tração, logo F = T1 = 240 N. É a alternativa (d), igual ao gabarito.%
\end{itemize}%
\paragraph{Pontos de dúvida}%
\begin{itemize}%
\item%
\textit{(solution)} A resolução fala em uma "normal N do piso" que se anularia na iminência. Na figura original, porém, o bloco já está suspenso, bem acima da base, sem tocar o piso. O resultado não muda, porque com o bloco pendurado e em repouso vale T = mg do mesmo jeito. Mas o argumento da normal descreve uma situação que a figura não mostra.%
\item%
\textit{(statement)} O enunciado diz "iminência de ser levantada", mas a figura mostra o bloco já pendurado. A leitura natural é a de equilíbrio estático (F mínima para sustentar ou começar a erguer o bloco), e nela F = 2mg = 240 N. Nenhuma outra interpretação razoável leva a outra alternativa.%
\end{itemize}%
\section*{Questão 66}%
\begin{longtable}{@{}p{0.06\textwidth} p{0.33\textwidth} p{0.33\textwidth} p{0.1\textwidth} p{0.08\textwidth}@{}}\toprule Item & Resolução & Gabarito & Veredito & Certo \\ \midrule
Q66 & (C) 4,0 N --- atrito cinético f\_at = $\mu$\_C$\cdot$m$\cdot$g = 0,02$\cdot$20$\cdot$10 = 4,0 N & C & igual & ambos \\
\bottomrule\end{longtable}%
\begin{itemize}%
\item%
\textbf{Q66)} Conferi a transcrição com a página 22 do PDF (página 20 da prova) e ela confere com o original. Com o disco em movimento, o atrito é cinético. Como a pista é horizontal, N = mg = 200 N, e então f\_at = 0,02$\cdot$200 = 4,0 N, que é a alternativa C. Os distratores também estão bem explicados: 0,4 N = $\mu$\_C$\cdot$m, 2,0 N = $\mu$\_E$\cdot$m e 20 N = $\mu$\_E$\cdot$mg, que é o atrito estático máximo.%
\end{itemize}%
\paragraph{Pontos de dúvida}%
\begin{itemize}%
\item%
\textit{(solution)} A resolução supõe, sem dizer, que a ação do esfregão aparece só no $\mu$\_C informado. Também supõe que nenhuma outra força vertical atua sobre o disco, de modo que N = mg. São hipóteses naturais para o enunciado e não mudam a resposta.%
\end{itemize}%
\section*{Questão 67}%
\begin{longtable}{@{}p{0.06\textwidth} p{0.33\textwidth} p{0.33\textwidth} p{0.1\textwidth} p{0.08\textwidth}@{}}\toprule Item & Resolução & Gabarito & Veredito & Certo \\ \midrule
Q67 & (B): patamar de fusão em 30 $^\circ$C e T3 = 50 $^\circ$C em t = 230 s & B & igual & ambos \\
\bottomrule\end{longtable}%
\begin{itemize}%
\item%
\textbf{Q67)} Conferi com a página 23 do PDF original. Os quatro gráficos correspondem à transcrição. Na fase líquida, Q = m$\cdot$c$\cdot$$\Delta$T, então $\Delta$T = 80/(10$\cdot$0,4) = 20 $^\circ$C e T3 = 50 $^\circ$C. (b) é o único gráfico que mostra aquecimento do sólido de 20 a 30 $^\circ$C (0 a 10 s), patamar de fusão a 30 $^\circ$C (até 210 s) e subida até 50 $^\circ$C em t $\approx$ 230 s. Os descartes da resolução estão certos: (a) termina em 120 $^\circ$C, o que exigiria Q = 360 J; (c) não tem patamar de fusão e chega a 50 $^\circ$C em $\approx$ 220 s; (d) mostra resfriamento. A observação de que a potência é variável, e por isso as inclinações não precisam ser comparadas, também está correta.%
\end{itemize}%
\paragraph{Pontos de dúvida}%
\begin{itemize}%
\item%
\textit{(statement)} O instante t = 230 s em que o gráfico (b) chega a 50 $^\circ$C foi lido no gráfico. O eixo tem marcas em 200, 220 e 240, e a reta termina entre 220 e 240, compatível com 230 s. Isso não muda a escolha, porque só (b) tem T3 = 50 $^\circ$C junto com o patamar de fusão.%
\end{itemize}%
\section*{Questão 68}%
\begin{longtable}{@{}p{0.06\textwidth} p{0.33\textwidth} p{0.33\textwidth} p{0.1\textwidth} p{0.08\textwidth}@{}}\toprule Item & Resolução & Gabarito & Veredito & Certo \\ \midrule
única & (B) 0,25 A & B & igual & ambos \\
\bottomrule\end{longtable}%
\begin{itemize}%
\item%
\textbf{única)} Conferi o enunciado na página 23 do PDF: P = 3,0 W, U = 12 V e as alternativas estão corretas na transcrição. Pela relação P = U$\cdot$i, que vale para qualquer dispositivo, i = 3,0/12 = 0,25 A, que é a alternativa B. A resolução tem a física e as unidades corretas (W/V = A), confere o resultado (12$\cdot$0,25 = 3,0 W) e explica bem por que as outras alternativas estão erradas: c) é U/P, d) é P$\cdot$U e a) dá 0,36 W.%
\end{itemize}%
\section*{Questão 69}%
\begin{longtable}{@{}p{0.06\textwidth} p{0.33\textwidth} p{0.33\textwidth} p{0.1\textwidth} p{0.08\textwidth}@{}}\toprule Item & Resolução & Gabarito & Veredito & Certo \\ \midrule
única (Q69) & (D) 30 mA & D & igual & ambos \\
\bottomrule\end{longtable}%
\begin{itemize}%
\item%
\textbf{única (Q69))} Reli o original (p. 21--22 do caderno, que são as páginas 23--24 do PDF). A transcrição está fiel. São quatro ramos em paralelo ligados diretamente aos 9 V, e a chave do ramo de R4 se fecha quando o alarme é ativado. R\_eq = 1,2 k$\Omega$/4 = 300 $\Omega$ e i = 9,0/300 = 30 mA. Conferi também os distratores, que correspondem ao que a resolução diz: (a) 9/4800 $\approx$ 1,875 mA, os quatro resistores em série; (b) 9/1200 = 7,5 mA, a corrente de um ramo só; (c) 9/400 = 22,5 mA, três ramos com a chave aberta.%
\end{itemize}%
\paragraph{Pontos de dúvida}%
\begin{itemize}%
\item%
\textit{(solution)} A resolução supõe bateria ideal, fios e chave de resistência desprezível. O enunciado não diz isso, mas é a hipótese usual e é a que leva ao gabarito.%
\end{itemize}%
\section*{Questão 70}%
\begin{longtable}{@{}p{0.06\textwidth} p{0.33\textwidth} p{0.33\textwidth} p{0.1\textwidth} p{0.08\textwidth}@{}}\toprule Item & Resolução & Gabarito & Veredito & Certo \\ \midrule
Questão 70 & D & D & igual & ambos \\
\bottomrule\end{longtable}%
\begin{itemize}%
\item%
\textbf{Questão 70)} A transcrição confere com a página 24 do PDF (página 22 do caderno). No trecho, Habermas diz que justiça e solidariedade "remontam à mesma raiz da moral", a vulnerabilidade dos seres vivos. Diz também que esses seres "só se particularizam em indivíduos por ação da socialização". Por isso a moral, ao defender os direitos do indivíduo, tem de defender também o bem{-}estar da comunidade. A justiça trata dos direitos iguais e da liberdade subjetiva dos indivíduos; a solidariedade trata das partes ligadas por uma forma de vida partilhada intersubjetivamente. Isso é exatamente o que diz a (d). As outras alternativas contradizem o texto. A (a) fala em contradição insolúvel, mas o texto mostra uma raiz comum. A (b) diz que a modernidade impede a vida ética, o que o texto não afirma. A (c) diz que a ética do discurso se opõe à vulnerabilidade e que a particularização independe da socialização, o contrário do texto. A análise de cada alternativa na resolução está correta.%
\end{itemize}%
\section*{Questão 71}%
\begin{longtable}{@{}p{0.06\textwidth} p{0.33\textwidth} p{0.33\textwidth} p{0.1\textwidth} p{0.08\textwidth}@{}}\toprule Item & Resolução & Gabarito & Veredito & Certo \\ \midrule
71 & (C) A força desumaniza aqueles sobre os quais recai e aqueles que dela se utilizam para transformar seres humanos em coisas. & C & igual & ambos \\
\bottomrule\end{longtable}%
\begin{itemize}%
\item%
\textbf{71)} Confiro o enunciado com a página 22 do PDF (página 24 do arquivo) e a transcrição está fiel. A alternativa C é a que se apoia no Texto 1. A força transforma em coisa quem é submetido a ela: vira cadáver, ou continua vivo, mas como "pedra". Ela também "embriaga aquele que a possui, ou julga possuí{-}la" e "ninguém a possui verdadeiramente", então atinge os dois lados. A está errada porque Weil nega que alguém detenha a força. B está errada porque Weil chama as palavras de Aquiles de "razoáveis", ou seja, lúcidas e humanas. D está errada porque bois e ovelhas aparecem como bens recuperáveis, em contraste com a vida humana, que não se recupera. O texto não estabelece uma hierarquia entre o homem e os animais. A análise de cada alternativa na resolução está correta.%
\end{itemize}%
\paragraph{Pontos de dúvida}%
\begin{itemize}%
\item%
\textit{(solution)} A resolução diz que quem exerce a força é desumanizado, porque fica "embriagado" e vira instrumento da força. Isso é uma inferência a partir de "embriaga aquele que a possui" e "ninguém a possui verdadeiramente". O Texto 1 não diz isso com essas palavras. É a leitura que o próprio gabarito exige na alternativa C, então a resposta não muda. Ainda assim, "vira seu instrumento" é interpretação do resolvedor, e não citação do texto.%
\end{itemize}%
\section*{Questão 72}%
\begin{longtable}{@{}p{0.06\textwidth} p{0.33\textwidth} p{0.33\textwidth} p{0.1\textwidth} p{0.08\textwidth}@{}}\toprule Item & Resolução & Gabarito & Veredito & Certo \\ \midrule
72 & (A) o ego moderno não se constitui apenas pelo Penso, logo existo, de Descartes, mas também pela conquista e colonização do Outro. & A & igual & ambos \\
\bottomrule\end{longtable}%
\begin{itemize}%
\item%
\textbf{72)} Conferi o enunciado na página 24 do PDF (numerada 22 na prova), e a transcrição está fiel. Segundo o texto, a Modernidade nasce quando a Europa se define "como um ego descobridor, conquistador, colonizador da Alteridade constitutiva da própria Modernidade". A alternativa (a) diz exatamente isso, contrastando com o cogito cartesiano, que é o marco tradicional do sujeito moderno. As outras estão erradas. Em (b), o diálogo é uma tarefa que ainda falta cumprir ("A tarefa, portanto, é desenvolver uma filosofia do diálogo"), e não algo que a filosofia moderna já tenha feito. Em (c), a modernidade nasce nas "cidades europeias medievais, livres, centros de enorme criatividade", e não "dos escombros da Idade Média". Em (d), a filosofia da libertação é a do oprimido e do excluído, não uma mediação entre europeus e colonizados para "dirimir" a violência. A análise da resolução está correta e bem apoiada no texto.%
\end{itemize}%
\paragraph{Pontos de dúvida}%
\begin{itemize}%
\item%
\textit{(solution)} A resolução usa um conhecimento que não está no texto: o conceito de 'ego conquiro', que viria antes do 'ego cogito' na obra de Dussel. A informação está correta, mas não é necessária para a questão. A alternativa (a) já se justifica só pela expressão 'ego descobridor, conquistador, colonizador' do trecho citado.%
\end{itemize}%
\end{document}